% concatenated from GConfR.tex
\documentclass[11pt]{amsart}
\usepackage[top = 1in, bottom = 1in, left = 1.5in, right=1.5in,
marginparwidth=1.3in]{geometry}
%\setlength{\parindent}{.5in}
\usepackage{setspace}
\usepackage{amsmath, amsfonts, amsthm, amstext, xspace}
\usepackage[mathscr]{euscript}
\usepackage{amssymb,latexsym}
\usepackage{comment}
\usepackage[dvipsnames,svgnames]{xcolor}
\usepackage{hyperref}
\usepackage[capitalise]{cleveref}
\usepackage{adjustbox}

%For diacritic in Piotr's name
\usepackage[utf8]{inputenc}
\usepackage[T1]{fontenc}


\definecolor{seagreen}{RGB}{46,139,87}
\definecolor{maroon}{RGB}{128,0,0}
\definecolor{darkviolet}{RGB}{148,0,211}
\definecolor{twelve}{RGB}{100,100,170}
\definecolor{thirteen}{RGB}{100,150,50}
\definecolor{fourteen}{RGB}{200,0,0}
\definecolor{fifteen}{RGB}{0,200,0}
\definecolor{sixteen}{RGB}{0,0,200}
\definecolor{seventeen}{RGB}{200,0,200}
\definecolor{eighteen}{RGB}{0,200,200}

\usepackage{centernot}
\usepackage{longtable}
\usepackage{mathtools}
\usepackage{stmaryrd}
\makeatletter
\newcommand{\xMapsto}[2][]{\ext@arrow 0599{\Mapstofill@}{#1}{#2}}
\def\Mapstofill@{\arrowfill@{\Mapstochar\Relbar}\Relbar\Rightarrow}
\makeatother

\usepackage[nolabel]{showlabels} % uncomment (and comment out the next line) to remove \label{} labels
%\usepackage[inline]{showlabels}
\renewcommand{\showlabelfont}{\tiny\tt\color{olive}}

% J.D. comments

\newcommand{\mjd}[1]{\marginpar{\color{red}#1}}
\newcommand{\cjd}[1]{\begin{quote}{\color{red}[\bf J.D.: #1]} \end{quote}}
\newcommand{\jd}[1]{{\color{red}{#1}}}

%\usepackage{xcolor}



\usepackage[all]{xy}
\newtheorem{thm}[equation]{Theorem}
\newtheorem*{theorem*}{Theorem}
\newtheorem*{conjecture*}{Conjecture}
\newtheorem*{corollary*}{Corollary}
\newtheorem{lemma}[equation]{Lemma}
\newtheorem{theorem}[equation]{Theorem}
\newtheorem{corollary}[equation]{Corollary}
\newtheorem{statement}[equation]{Statement}
\newtheorem{proposition}[equation]{Proposition}
\newtheorem{conjecture}[equation]{Conjecture}
\newtheorem{hypothesis}[equation]{Hypothesis}
\newtheorem{convention}[equation]{Convention}
\newtheorem{assumption}[equation]{Assumption}
\theoremstyle{definition}
\newtheorem{definition}[equation]{Definition}
\newtheorem{example}[equation]{Example}
\newtheorem{remark}[equation]{Remark}
\newtheorem{procedure}[equation]{Procedure}
\newtheorem{construction}[equation]{Construction}
\newtheorem{question}[equation]{Question}
\newcommand{\ra}{\rightarrow}
\newcommand{\la}{\leftarrow} 
\newcommand{\lra}{\longrightarrow}
\newcommand{\lla }{\longrightarrow}
\numberwithin{equation}{subsection} % Number equations by section, like (1.1) instead of (1)

\newtheorem{thmx}{Theorem}
\renewcommand{\thethmx}{\Alph{thmx}} 


%PUT YOUR MACROS HERE

%Alphabet shortcuts
\def\a{\mathbb{A}}
\def\b{\mathbb{B}}
\def\c{\mathbb{C}}
\def\e{\mathbb{E}}
\def\f{\mathbb{F}}
\def\g{\mathbb{G}}
\def\h{\mathbb{H}}
\def\H{\mathcal{H}}
%\def\k{\mathbb{K}}
\def\l{\mathbb{L}}
\def\m{\mathbb{M}}
\def\n{\mathbb{N}}
\def\p{\mathbb{P}}
\def\q{\mathbb{Q}}
\def\r{\mathbb{R}}
\def\s{\mathbb{S}}
\def\t{\mathbb{T}}
\def\z{\mathbb{Z}}
\def\N{\mathbb{N}}

\def\ca{\mathcal{A}}
\def\cb{\mathcal{B}}
\def\cc{\mathcal{C}}
\def\cd{\mathcal{D}}
\def\ce{\mathcal{E}}
\def\cf{\mathcal{F}}
\def\cg{\mathcal{G}}
\def\cm{\mathcal{M}}
\def\cn{\mathcal{N}}
\def\co{\mathcal{O}}
\def\cp{\mathcal{P}}
\def\cq{\mathcal{Q}}
%\def\cr{\mathcal{R}}
\def\cs{\mathcal{S}}
\def\ct{\mathcal{T}}
\def\cu{\mathcal{U}}
\def\cv{\mathcal{V}}

\def\sp{\mathscr{P}}

\def\fa{\mathfrak{a}}
\def\fb{\mathfrak{b}}
\def\fg{\mathfrak{g}}
\def\fh{\mathfrak{h}}
\def\fm{\mathfrak{m}}
\def\fp{\mathfrak{p}}
\def\fq{\mathfrak{q}}
\def\fu{\mathfrak{u}}

\def\uI{\underline{I}}
\def\uM{\underline{M}}
\def\uP{\underline{P}}
\def\uR{\underline{R}}


%Commonly used terms
\def\zmodp{\mathbb{Z}/p}
\def\zlocp{\mathbb{Z}_{(p)}}
\def\ctp{\widehat{\otimes}}
\def\uA{\underline{A}}
\def\uC{\underline{C}}
\def\uH{\underline{H}}
\def\uM{\underline{M}}

%category abbreviations
\def\Moda{\mathcal{M}od_A}
\def\Spt{\mathcal{S}pt}
\def\Motc{\mathcal{M}ot_\c}
\def\Motr{\mathcal{M}ot_\r}
\def\Sptc{\mathcal{S}pt^{C_2}}
\def\Ab{\mathcal{A}b}
\def\mft{\underline{\mathbb{F}_2}}

%homotopy theory symbols
\def\holim{\operatorname{lim}}
\def\Gal{\operatorname{Gal}}
\def\hocolim{\operatorname{colim}}
\def\THH{\operatorname{THH}}
\def\TC{\operatorname{TC}}
\def\HH{\operatorname{HH}}
\def\HC{\operatorname{HC}}
\def\Ho{\operatorname{Ho}}
\def\Spec{\operatorname{Spec}}
\def\Spf{\operatorname{Spf}}
\def\chara{\operatorname{char}}
\def\Tor{\operatorname{Tor}}
\def\Ext{\operatorname{Ext}}
\def\Hom{\operatorname{Hom}}
\def\colim{\operatorname{colim}}
\def\cofib{\operatorname{cofib}}
\def\fib{\operatorname{fib}}
\def\ker{\operatorname{ker}}
\def\coker{\operatorname{coker}}
\def\id{\operatorname{id}}
\def\re{\operatorname{Re}}
\def\SH{\operatorname{SH}}
\def\Sm{\operatorname{Sm}}
\def\eff{\operatorname{eff}}
\def\veff{\operatorname{veff}}
\def\sc{\operatorname{sc}}
\def\vsc{\operatorname{vsc}}
\def\gr{\operatorname{gr}}
\def\mASS{\operatorname{mASS}}
\def\mANSS{\operatorname{mANSS}}
\def\eSSS{\operatorname{eSSS}}
\def\aESSS{\operatorname{aESSS}}
\def\aVSSS{\operatorname{aVSSS}}
\def\im{\operatorname{im}}
\def\cl{\operatorname{cl}}
\renewcommand{\Re}{\operatorname{Re}}
\def\Ad{\operatorname{Ad}}
\def\mfp{\underline{\mathbb{F}_p}}
\def\mfz{\underline{\mathbb{Z}}}
\newcommand{\Conf}{\operatorname{Conf}}
\newcommand{\Aut}{\operatorname{Aut}}
\def\Sp{\operatorname{Sp}}
\def\upi{\underline{\pi}}
\def\Emb{\operatorname{Emb}}
\def\tr{\operatorname{tr}}
\def\res{\operatorname{res}}

%\newcommand{\limt}[1]{\lim_{#1 \in \mathbb{N}^{\text{op}}}}
\newcommand{\limt}[1]{\underset{#1}{\lim \,}}
\def\vbar{\overline{v}}

% Eva's newcommands
\usepackage{marginnote}
\usepackage{soul}
\newcommand{\sm}{\wedge}
\newcommand{\iy}{\infty}
\newcommand{\attop}[1]{{\let\textstyle\scriptstyle\let\scriptstyle\scriptscriptstyle\substack{#1}}}
\renewcommand{\atop}[1]{{\let\scriptstyle\textstyle\let\scriptscriptstyle\scriptstyle\substack{#1}}}
\newcommand{\x}{\times}
\renewcommand{\st}{\hspace{2pt} : \hspace{2pt}}
\newcommand{\proj}{\operatorname{proj}}
\newcommand{\contains}{\supseteq}
\renewcommand{\l}{\overset}
\newcommand{\too}[1]{\l{#1}\to}
\newcommand{\hteq}{\simeq}
\newcommand{\ints}{\cap}
\newcommand{\isom}{\cong}
\newcommand{\Z}{\mathbb{Z}}
\newcommand{\F}{\mathbb{F}}
\newcommand{\hocofib}{\operatorname{hocofib}}
\newcommand{\til}[1]{\widetilde{#1}}
\newcommand{\supp}{\operatorname{supp}}


\makeatletter
\newcommand{\switchmargin}{
\if@reversemargin
\normalmarginpar
\else
\reversemarginpar
\fi
}
\makeatother
\newcommand{\eva}[1]{{\color{teal}#1}}
\newcommand{\EvaNote}[1]
   {\marginpar{\raggedright\smaller \smaller \color{teal}#1}\switchmargin}
\newcommand{\highlighteva}[1]{\ifmmode{\text{\sethlcolor{llgray}\hl{$#1$}}}\else{\sethlcolor{llteal}\hl{#1}}\fi}
\newcommand{\EvaNoteHl}[2]{\marginnote{\smaller \smaller \color{teal}
#2}{\highlighteva{#1}}\switchmargin}



\definecolor{llteal}{RGB}{198,232,227}
\definecolor{llred}{RGB}{237,228,228}
\definecolor{llgray}{RGB}{230,230,230}
\definecolor{maroon}{RGB}{150,0,0}
\definecolor{orange}{RGB}{255,165,0}
\newcommand{\edit}[1]{{\it{\color{gray}#1}}}
\newcommand{\fixme}[1]{{\color{maroon}\it{#1}}}
\newcommand{\citeme}[2][\!\!]{{\color{orange}[#2~\textit{#1}]}}
\newcommand{\highlight}[1]{\ifmmode{\text{\sethlcolor{llgray}\hl{$#1$}}}\else{\sethlcolor{llred}\hl{#1}}\fi}
\newcommand{\fixmehl}[2]{\marginnote{\smaller \smaller\color{maroon} #2}{\highlight{#1}}}

% Chase's comments & commands

\newcommand{\chase}[1]{{\color{purple}#1}}
\newcommand{\ChaseNote}[1]
   {\marginpar{\raggedright\smaller \smaller \color{purple}#1}\switchmargin}

\newcommand{\ZZ}{\mathbb Z}
\newcommand{\sma}{\wedge}
\newcommand{\sus}{\Sigma}
\newcommand{\nv}{^{-1}}

\usepackage{amssymb}
\usepackage{tikz-cd}

\setcounter{tocdepth}{1}

\author{Eva Belmont}\address{Case Western Reserve University}\email{eva.belmont@case.edu}
\author{J.D. Quigley}\address{University of Virginia}\email{mbp6pj@virginia.edu}
\author{Chase Vogeli}\address{Cornell University}\email{cpv29@cornell.edu}
\title{Bredon homological stability for configuration spaces of $G$-manifolds}

\begin{document}

\begin{abstract}
McDuff and Segal proved that unordered configuration spaces of open manifolds satisfy
homological stability: there is a stabilization map $\sigma: C_n(M)\to
C_{n+1}(M)$ which is an isomorphism on $H_d(-;\Z)$ for $n\gg d$.
For a finite group $G$ and an open $G$-manifold $M$, under some hypotheses we
define a family of equivariant stabilization maps $\sigma_{G/H}:C_n(M)\to
C_{n+|G/H|}(M)$ for $H\leq G$. In general, these do not induce
stability for Bredon homology, the equivariant analogue of singular homology. Instead, we show that each $\sigma_{G/H}$ induces
isomorphisms on the ordinary homology of the fixed points of $C_n(M)$,
and if the group is Dedekind (e.g. abelian), we obtain the following Bredon
homological stability statement: $H^G_d(\bigsqcup_{n\geq 0}C_n(M))$ is finitely
generated over $\Z[\sigma_{G/H} : H\leq G]$. This reduces to the classical
statement when $G=e$.
\end{abstract}

\maketitle

\tableofcontents

\section{Introduction}

Let $X$ be a topological space and let $C_k(X)$ denote the configuration space of $k$ unordered points in $X$. If $X$ is the interior of a manifold with nonempty boundary, then there are stabilization maps
\begin{equation}\label{Eqn:StabIntro}
\sigma: C_k(X) \to C_{k+1}(X)
\end{equation}
defined by ``adding a point near the boundary" (cf. \cref{SS:Conf}). In the 1970's, McDuff \cite{McD75} and Segal \cite{Seg79} observed that these stabilization maps induce isomorphisms in a range of integral homology groups:

\begin{thm}[McDuff--Segal, Stability for Configuration Spaces (Strong Form)]\label{Thm:MS}
Assume $M$ is the interior of a manifold with nonempty boundary. The stabilization map \eqref{Eqn:StabIntro} induces an isomorphism
$$\sigma_*: H_d (C_k(M)) \to H_d(C_{k+1}(M))$$
for $d \leq k/2$. 
\end{thm}


\subsection{Statement of main results}

In this paper, we investigate an equivariant analogue of the McDuff--Segal theorem where the manifold $M$ is replaced by a $G$-manifold (where $G$ is a finite abelian group) and singular homology is replaced by Bredon homology. 

The Bredon homology of a $G$-space $X$ with coefficients in the constant Mackey functor $\mfz$, written $H^G_*(X; \mfz)$, is an equivariant analogue of singular homology with integer coefficients. We recall this definition of Bredon homology, as well as its representation by a $G$-spectrum, in \cref{SS:Bredon}.

If $X$ is a $G$-space, then the space of unordered configurations of $k$ points in $X$
$$C_k(X) := \{(x_1,\ldots,x_k) \in X^k: x_i \neq x_j \text{ for all } i\neq j\}/\Sigma_k$$
is also a $G$-space with $G$ acting diagonally. If $X$ is a $G$-manifold which is the interior of a $G$-manifold with nonempty boundary containing a fixed point, then we can define $G$-equivariant stabilization maps
\begin{equation}\label{Eqn:FPGIntro}
\sigma_{G/G}: C_k(X) \to C_{k+1}(X)
\end{equation}
by ``adding a fixed point near the boundary'' (see \cref{SS:GConf}). In this situation, we say $X$ is \emph{$G$-stabilizable} (\cref{Def:HStab}). We can then ask if these maps induce isomorphisms in Bredon homology for $k$ sufficiently large. Surprisingly, this is not the case:

\begin{thmx}[Failure of Strong Equivariant Form of McDuff--Segal, Proposition \ref{Prop:H0Unstable}]\label{MT:BredonFails}
Let $G=C_p$ and let $M$ be a $G$-manifold which is $G$-stabilizable. The map
\begin{equation}\label{Eqn:GStabIntro}
(\sigma_{G/G})_*: H_0^G(C_n(M); \mfz) \to H_0^G(C_{n+1}(M); \mfz)
\end{equation}
is not surjective for any $n \geq 1$.
\end{thmx}

Although the strong form of the McDuff--Segal Theorem fails equivariantly, there is a weaker form which we can extend to the equivariant setting. Let
$$C(X) := \bigsqcup_{k \geq 0} C_k(X)$$
denote the space of all unordered configurations of finitely many points in $X$ equipped with the disjoint union topology. The nonequivariant stabilization maps \eqref{Eqn:StabIntro} give rise to a self-map
\begin{equation}\label{Eqn:BigStabIntro}
\sigma: C(X) \to C(X)
\end{equation}
which allows us to state the following corollary of the McDuff--Segal Theorem:

\begin{thm}[McDuff--Segal, Nonequivariant Stability (Finite Generation Form)]
Assume $M$ is an open, connected manifold with $\dim(M)\geq2$. Moreover, suppose that $H_d(M)$ is finitely generated as an abelian group for all $d\geq 0$. Then the homology group $H_d(C(M))$ is finitely generated as a $\z[\sigma_*]$-module for all $d\geq 0$. 
\end{thm}

\begin{remark}
The assumption that $H_d(M)$ is finitely generated ensures that $H_d(C_k(M))$ is finitely generated for all $k \geq 0$ (\cref{Prop:MtoCnM}), so the finite generation form of nonequivariant stability is equivalent (cf. \cref{Prop:AlgStability}) to saying that the map \eqref{Eqn:StabIntro} is an isomorphism for all $k \gg 0$. 

We note that the philosophy of encoding stability as a finite generation condition is not new. It has become increasingly prevalent in the literature on homological stability (e.g., \cite[Theorem 1.1]{ADCK20}) and representation stability (e.g. \cite[Theorem 6.2.1]{CEF-FI}).
\end{remark}

The obvious finite generation form of equivariant stability fails since the equivariant stabilization map \eqref{Eqn:GStabIntro} does not induce isomorphisms for all $k \gg 0$. In other words, $H_d^G(C(M); \mfz)$ is \emph{not} finitely generated as a $\z[(\sigma_{G/G})_*]$-module. 

However, there is a naturally occurring polynomial ring over which $H_d^G(C(M); \mfz)$ is a finitely generated module. If $X$ is a $G$-manifold which is the interior of a $G$-manifold with boundary containing an orbit of the form $G/H$, then we can define a new equivariant stabilization map
\begin{equation}\label{Eqn:GHStabIntro}
\sigma_{G/H}: C_k(M) \to C_{k+|G/H|}(M)
\end{equation}
by ``adding an \emph{orbit} of type $G/H$ near the boundary.'' In this situation, we say $X$ is \emph{$H$-stabilizable} (\cref{Def:HStab}). 

We adopt the philosophy that since orbits are equivariant analogues of points, the equivariant analogue of stabilization maps should be the family of stabilization maps $\sigma_{G/H}$. As such, we replace the ring $\z[\sigma_*]$ by the ring 
$$P_G=\z[(\sigma_{G/H})_*: (H) \leq G]$$
of all equivariant stabilization maps. By considering all possible equivariant stabilization maps, we obtain our stability result.

%In our reduction of equivariant homological stability for $C(M)$ to classical homology stability, we must consider the spaces $M_{(H)}$ of spaces with isotropy conjugate to $H$. See \cref{Cor:CGSDecomp} for how these spaces arise.

\begin{thmx}[{Bredon Homological Stability, \cref{thm:B}}]\label{MT:Stability}
Let $G$ be an abelian group and $M$ a $G$-manifold that is the interior of a
compact $G$-manifold with boundary. Assume that for all $H\leq G$, $M$ is
$H$-stabilizable and $M^H$ is connected.
Then $H_d^H(C(M);\underline{\Z})$ is finitely generated over $P_G$ for all $H\leq G$.
\end{thmx}
For example, these hypotheses are satisfied in the case where $M$ is a sum of
copies of the regular representation (see Example \ref{ex:rho}).

\begin{remark}
See Lemma \ref{lem:RK-to-RG} for the main reason for the abelian hypothesis, and
Remark \ref{rmk:abelian} for a mild relaxation of this hypothesis.
Group-theoretic hypotheses are also used in Lemma
\ref{lem:hypothesis-translation} and \ref{lem:fg-homology}.
%\cref{lem:hypothesis-translation}
%discusses an alternate condition which implies the connectedness hypothesis about
%$M_{(H)}/G$.
\end{remark}

\begin{remark}
We discuss two natural extensions of \cref{MT:Stability} in \cref{Sec:Gen}: in \cref{Thm:ROG}, we show that stability holds for $RO(G)$-graded Bredon homology, and in \cref{Thm:MF}, we show that stability holds for Mackey functor-valued Bredon homology. The $RO(G)$-graded Bredon homology of \emph{ordered} configuration spaces of $G$-manifolds will also be studied in forthcoming work of Dugger and Hazel \cite{DH23}. 
\end{remark}

%For a $G$-space $X$, one natural notion of homology is \emph{$\z$-graded Bredon homology}
%$$H^G_*(X; \uM)$$
%with coefficients in a Mackey functor $\uM$. Roughly speaking, a Mackey functor $\uM$ assigns an abelian group $\uM(H)$ to each subgroup $H \leq G$, along with some additional structure relating these assignments for subconjugate subgroups. The Bredon homology of $X$ is itself a Mackey functor; its value on the trivial subgroup
%$$H^G_*(X;\uM)(e) \cong H_*(X;\uM(e))$$
%coincides with ordinary singular homology with coefficients in (the abelian group) $\uM(e)$. If $M$ is an open $G$-manifold where $\sigma$ is defined and $\uM$ is a Mackey functor with $\uM(e)=\z$, then the map
%$$\sigma_* : H_i^G(C_k(M);\uM)(e) \cong H_i(C_k(M); \z) \to H_i(C_{k+1}(M); \z) \cong H_i^G(C_{k+1}(M); \uM)(e)$$
%is an isomorphism for $i \leq k/2$ by the McDuff--Segal theorem.

%However, we can ask for a refinement. The value of the Bredon homology $H_*^G(X;\uM)$ on a nontrivial subgroup $e \neq H \leq G$ is related to the homology of the fixed points $X^H$, but it is in general more complicated than just $H_*(X^H; \uM(H))$. We explore the following extension of classical homological stability:

%\begin{question}
%Let $M$ be an open $G$-manifold. Is there a nontrivial Mackey functor $\uM$ and a function $N(k)$ depending on $k$ for which the map
%$$\sigma_* : H_i^G(C_k(M); \uM) \to H_i^G(C_{k+1}(M); \uM)$$
%induces an isomorphism \emph{of Mackey functors}? 
%\end{question}

%Surprisingly, we find that this strengthening of homological stability fails:

%\begin{thmx}\label{MT:BredonFails}
%Let $G$ be a finite group with nontrivial Sylow $p$-subgroups for $p \in P$, let $M$ be an open $G$-manifold, and let $\uM$ be any $G$-Mackey functor. The map
%$$\sigma_*: H_i^G(C_k(M); \uM) \to H_i^G(C_{k+1}(M); \uM)$$
%is injective for all $i \geq 0$, but fails to be surjective whenever $p|(k-1)$ for some $p \in P$. 
%\end{thmx}

\subsection{Motivation}

Before discussing the proofs of \cref{MT:BredonFails} and \cref{MT:Stability}, we discuss some motivating ideas and directions for future work.

\subsubsection{Equivariant loop spaces and Dyer--Lashof operations}

Nonequivariantly, May's recognition principle \cite{May72} implies that $n$-fold loop spaces are equivalent to algebras over the $E_n$-operad $\ce_n$ whose $k$-th space is $C_k(\r^n)$. Using Cohen's computation of $H_*(C_k(\r^n))$ in \cite[Ch. III]{CLM76}, one can then define Dyer--Lashof operations on the mod $p$ homology of any $n$-fold loop space. 

In \cite{GM17}, Guillou and May proved an equivariant analogue of the recognition principle: if $V$ is a real orthogonal $G$-representation, then $V$-fold loop spaces are equivalent to algebras over the $E_V$-operad $\ce_V$ whose $k$-th space is $C_k(V)$. They remark that in contrast with the nonequivariant situation, very little is known about $H_d^G(C_k(V))$. \cref{MT:Stability} gives the first step toward a more systematic analysis of these Bredon homology groups which could shed light on equivariant Dyer--Lashof operations (cf. \cite{Wil19} for $G=C_2$). We verify that the hypotheses of \cref{MT:Stability} are satisfied in the case of sums of regular representations in \cref{ex:rho}.


\subsubsection{Equivariant diffeomorphism groups of smooth $G$-manifolds}

Configuration spaces of ordinary manifolds appear in the study of diffeomorphism groups of smooth manifolds via the ``decoupling theorems" of B{\"o}digheimer--Tillmann \cite{BT01} and Bonatto \cite{Bon22}. It would be interesting to see if the configuration spaces of $G$-manifolds arise analogously in the equivariant diffeomorphism groups of smooth $G$-manifolds. 

% \subsubsection{Connections to mathematical physics}

% Configuration spaces of $G$-manifolds also arise naturally in mathematical physics. Recall the ``electrostatic map"
% $$C_k(\r^n) \to \Omega^n S^n$$
% which sends each configuration of $k$ points (regarded as particles with unit positive charge) to the electric field it generates \cite{Seg73}. The fact that this extends to a well-defined, equivariant electrostatic map (cf. \cite{RS00})
% $$C_k(V) \to \Omega^V S^V$$
% for any real orthogonal $G$-representation $V$ captures the symmetries of the electric field. More generally, if $M$ is a $G$-manifold, the observables on $C_k(M)$ which are equivariant with respect to the $G$-action on $C_k(M)$ are those which are suitably symmetric. 

%\subsubsection{Stability for the Amitsur--Dress homology of symmetric groups}

%\ChaseNote{I imagine this section will need to be rewritten}
%Another consequence of the McDuff--Segal Theorem is Nakaoka's celebrated result on the stable group homology of symmetric groups from \cite{Nak60}. As with configuration spaces, the inclusion $\Sigma_k \to \Sigma_{k+1}$ induces a map on homology 
%$$\sigma_*: H_*(\bigsqcup_{k \geq 0} \Sigma_k; \z) \to H_*(\bigsqcup_{k \geq 0} \Sigma_k; \z).$$
%The weak form of homological stability for unordered configuration spaces, along with the facts that
%$$C_k(\r^\infty) \simeq B\Sigma_k$$
%and
%$$H_d(\Sigma_k; \z) \cong H_d(B\Sigma_k; \z),$$
%recovers the following:

%\begin{thm}[Nakaoka]
%The $\z[\sigma_*]$-module $\bigoplus_{k \geq 0} H_d(B\Sigma_k; \z)$ (or equivalently, $\bigoplus_{k \geq 0} H_d(\Sigma_k; \z)$) is finitely generated for all integers $d \geq 0$.
%\end{thm}

%In equivariant homotopy theory, the classifying space $B\Sigma_k$ is replaced by the $G$-equivariant classifying space $B_G\Sigma_k$ (cf. \cref{Def:BGSigma}). There is a $G$-equivariant equivalence
%$$C_k(\rho^\infty) \simeq B_G \Sigma_k$$
%(cf. \cref{prop:BGSigma_n-configuration}) between configurations of $k$ unordered points in infinitely many copies of the regular representation of $G$ and $B_G\Sigma_k$, the $G$-equivariant classifying space for $\Sigma_k$.

%\begin{thmx}[{Equivariant Nakaoka Stability, \cref{Thm:Nakaoka}}]\label{MT:Nakaoka}
%The sum of Bredon homology groups $\bigoplus_{k \geq 0} H_d^G(B_G\Sigma_k; \mfz)$ is finitely generated as a $P_G$-module.  
%\end{thmx}

%\begin{remark}
%In equivariant algebra, group homology is replaced by Amitsur--Dress homology \cite{GM95}. \cref{MT:Stability} has a purely algebraic reformulation in terms of Amitsur--Dress homology via the isomorphism
%$$H_d^G(\Sigma; \mfz) \cong H_d^G(B_G\Sigma; \mfz)$$
%between the $G$-equivariant Amitsur--Dress homology of $\Sigma_k$ and the Bredon homology of $B_G\Sigma$.
%\end{remark}

%\begin{remark}
%The classical Barratt--Priddy Theorem \cite{BP72} gives an algebra isomorphism
%$$\colim_{k \to \infty} H_*(\Sigma_k; \z) \to \colim_{k \to \infty} H_*((\Omega^k \Sigma^k S^0)_0; \z)$$
%between the stable homology of symmetric groups and the homology of the space whose homotopy groups are the stable homotopy groups of spheres. This can be used to prove results about the stable homotopy groups of spheres, such as Serre's Finiteness Theorem \cite{Ser53}, via group homology. It could be interesting to study the applications of \cref{MT:Nakaoka} in conjunction with various equivariant analogues of the Barratt--Priddy Theorem \cite{BGS20, GM17}. 
%\end{remark}

%%Our motivation for studying the configuration spaces of $G$-spaces comes from equivariant homotopy theory. Nonequivariantly, May's recognition principle \cite{May72} implies that $n$-fold loop spaces are equivalent to algebras over the $E_n$-operad $\ce_n$ whose $k$-th space is $C_k(\r^n)$. Using this structure, it is possible to produce \emph{Dyer--Lashof operations}
%%$$Q_i : H_m(X;\f_2) \to H_{pm+i}(X;\f_2)$$
%%for $0 \leq i \leq n-1$ for any $n$-fold loop space $X$ (or similar operations in mod $p$ homology). The construction and analysis of these operations depends on Cohen's \cite[Ch. III]{CLM76} complete computation of the homology of $C_k(\r^n)$. 

%%\begin{remark}
%%Cohen's calculations can also be used to refine the McDuff--Segal stability theorem after inverting $2$, cf. \cite{KM15}. 
%%\end{remark}

%%In \cite{GM17}, Guillou and May proved an equivariant analogue of the recognition principle. For any real orthogonal $G$-representation $V$, we may form a ``representation sphere" $S^V$ by taking its one-point compactification, and so it makes sense to ask about $V$-fold loop spaces. Guillou and May showed that $V$-fold loop spaces are equivalent to algebras over the $E_V$-operad $\ce_V$ whose $k$-th space is $C_k(V)$. However, they note that unlike in the nonequivariant setting, very little is known about the Bredon homology of these equivariant configuration spaces. While a complete calculation along the lines of Cohen's work is beyond the scope of our paper, our results can be viewed as a step towards understanding Dyer--Lashof operations in Bredon homology. 

%%\begin{remark}
%%When $G = C_2$, Behrens and Wilson produced some Dyer--Lashof operations for Bredon homology in \cite{BW18}. We are unaware of similar constructions for other groups. 
%%\end{remark}

%%Although homological stability fails for the Bredon homology of configuration spaces of $G$-manifolds, there is more refined notion of equivariant configuration space where we \emph{do} obtain stability. In fact, we use these more refined configuration spaces to obtain a simple conceptual proof of \cref{MT:BredonFails}. 

\subsection{Overview of key ideas}

We now summarize the key ideas going into the proofs of \cref{MT:BredonFails} and \cref{MT:Stability}. 

\subsubsection{Equivariant stabilization maps}

A map between $G$-spaces induces a map on Bredon homology only if it is $G$-equivariant, so we must be careful about how we add points. One natural option is to add a \emph{fixed} point near infinity (provided our space has such a fixed point), which yields the equivariant stabilization map \eqref{Eqn:FPGIntro} mentioned above with underlying map \eqref{Eqn:StabIntro}. However, \cref{MT:BredonFails} shows that this map does not produce isomorphisms in Bredon homology after iteration. 

One of our key observations is that under suitable hypotheses on the manifold $M$, there is a stabilization map \eqref{Eqn:GStabIntro} for each subgroup $H \leq G$. \cref{MT:Stability} says that the Bredon homology stabilizes when all of these equivariant stabilization maps are taken into account simultaneously. 

\begin{remark}
The stabilization map corresponding to $H \leq G$ inserts an orbit of the form $G/H$, i.e., inserts $|G/H|$ points at once. Other stabilization maps which add multiple points at once have appeared in the context of higher order representation stability \cite{MW19}.
\end{remark}

\subsubsection{Formulating stability in terms of finite generation}

While none of the orbit stabilization maps \eqref{Eqn:GStabIntro} yield Bredon homological stability for a fixed $H$ (\cref{MT:BredonFails}), the collection of maps $\{ \sigma_{G/H} : (H) \leq G\}$ yield a form of stability (\cref{MT:Stability}). To make this precise, we translate homological stability from a statement about maps in homology to a statement about finite generation over a polynomial ring (\cref{Prop:AlgStability}). 

\subsubsection{Reduction from Bredon homology to singular homology of fixed points}
The analogue of cellular homology in the context of $G$-spaces is \emph{Bredon homology},
$H^G_d(-)$ (cf. \cite[XIII.4.1]{May96}, Section \ref{SS:Bredon}). The Bredon
homology of a $G$-space is not simply the homology of its fixed points, but in
Section \ref{Sec:Eqvt} we use equivariant stable homotopy theory to relate these
notions. The main lemma enabling this comparison is
\cref{Lemma:Bredon-to-PhiG}, together with statements in Section \ref{SS:GEM} relating the geometric localization of Bredon homology to singular homology.
% the (geometric) fixed points of equivariant Eilenberg-Maclane
%spaces to ordinary Eilenberg-Maclane spaces.

\subsubsection{Expressing fixed points in terms of nonequivariant configuration spaces}

We prove that the fixed points $C_n(M)^G$ satisfy homological stability by reducing to \cref{Thm:MS}. To do so, we express $C_n(M)^G$ in terms of configuration spaces of other manifolds using the notion of \emph{unordered $S$-configurations}. We note that we are not the first to consider equivariant configuration spaces (cf. \cite{RS00}), but we use a new convention here.

Let $X$ be a $G$-space and let $S$ be a finite $G$-set. We define the \emph{space of unordered $S$-configurations in $X$} (\cref{Def:CGS}) by
$$C_S^G(X) := \Emb^G(S,X)/\Aut(S).$$
When $M$ is $H$-stabilizable, the equivariant stabilization map \eqref{Eqn:GHStabIntro} restricts to a stabilization map 
$$\sigma_{G/H} : C_S^G(M) \to C_{S + G/H}^G(M),$$
where $S + G/H$ denotes the disjoint union of $S$ and $G/H$ as $G$-sets.

\begin{example}
If $S$ is a trivial $G$-set of cardinality $k$, then we have that
$$C_{S}^G(X) = C_k(X^G)$$
is the space of unordered configurations of $k$ points in the $G$-fixed points of $X$. If $X$ is a $G$-stabilizable $G$-manifold, the stabilization map
$$\sigma_{G/G} : C_{S}^G(X) \to C_{S + [G/G]}^G(X)$$
is the classical stabilization map
$$\sigma: C_k(X^G) \to C_{k+1}(X^G).$$
\end{example}

Passing through the $S$-configuration spaces $C_S^G(X)$ allows us to identify the fixed points $C_n(X)^G$ in terms of ordinary configuration spaces:

\begin{thmx}[{\cref{Prop:FixptConfDecomp} and \cref{Cor:CGSDecomp}}]\label{MT:CtoCSG}
\
\begin{enumerate}
\item For a $G$-space $X$ and $n\geq 0$, there is a disjoint union decomposition
$$ C_n(X)^G \cong \coprod_{|S|=n} C^G_S(X), $$
where the disjoint union runs over isomorphism classes of finite $G$-sets $S$ with cardinality $n$.

\item Let $S$ be the finite $G$-set
$$S = \bigsqcup_{(H)} k_{(H)}[G/H],$$
where $k_{(H)} \geq 0$ for each conjugacy class of subgroups $(H)$. Then
$$C_S^G(X) \cong \prod_{(H)} C_{k_{(H)}}(X_{(H)}/G),$$
where $X_{(H)}$ is the subspace of $X$ consisting of points with stabilizer conjugate to $(H)$ (\cref{Def:X_(H)}). 
\end{enumerate}
\end{thmx}

We apply part (1) of \cref{MT:CtoCSG} to prove \cref{MT:BredonFails}; the essential point is that when $G$ is nontrivial, the number of summands appearing in $C_n(X)^G$ increases as $n$ increases, so the stabilization map cannot be surjective on homology. We apply parts (1) and (2) to prove \cref{MT:CSGStability}.

\subsubsection{Stability for $S$-configurations and Bredon homology}

Since they can be expressed in terms of ordinary configuration spaces, the spaces of $S$-configurations often satisfy an analogue of the McDuff--Segal theorem:

\begin{thmx}[{\cref{Thm:CGSHomStab}}]\label{MT:CSGStability}
Let $M$ be a $G$-manifold and $S$ a finite $G$-set containing $k$
copies of the orbit
$[G/H]$ for a conjugacy class of subgroups $(H)$. If $M$ is $H$-stabilizable and
$M_{(H)}/G$ is connected, then the $(H)$-stabilization map of \eqref{eq:sigmaGH} induces an isomorphism
$$ (\sigma_{G/H})_*: H_d(C^G_S(M)) \to H_d(C^G_{S+[G/H]}(M)) $$
in integral homology in degrees $d\leq k/2$.
\end{thmx}

As explained above, \cref{MT:Stability} follows from stability for the fixed points $C_n(M)^G$. \cref{MT:CtoCSG} allows us to express these fixed points in terms of $S$-configurations, and \cref{MT:CSGStability} states that the equivariant stabilization map $\sigma_{G/H}$ from \eqref{Eqn:GHStabIntro} induces an isomorphism in the homology of $C_S^G(M)$ when $S$ contains sufficiently many copies of $[G/H]$. In particular, once $S$ contains sufficiently many copies of $[G/K]$ for \emph{any} conjugacy class of subgroups $(K)$, there is some equivariant stabilization map which induces an isomorphism on homology. Roughly speaking, \cref{MT:Stability} then follows from the fact that there are finitely many finite $G$-sets $S$ for which none of the equivariant stabilization maps induce isomorphisms in homology. 

\subsection{Outline}

In \cref{Sec:Eqvt}, we recall the equivariant homology theory (Bredon homology) we work with throughout the paper. We also recall the isotropy separation sequence and prove some key results (\cref{Lemma:Bredon-to-PhiG}, \cref{Prop:Green}) needed to reduce from the equivariant to nonequivariant setting. 

In \cref{Sec:Conf}, we recall configuration spaces, stabilization maps, and their equivariant analogues. We then introduce the more refined spaces of $S$-configurations and use them to prove \cref{MT:CtoCSG} which identifies the fixed points of the configuration spaces of $G$-manifolds with (sums of products of) configuration spaces in nonequivariant manifolds. 

In \cref{Sec:Stab}, we prove \cref{MT:BredonFails},  \cref{MT:Stability}, and \cref{MT:CSGStability}, using the results of the previous section to reduce to classical homological stability results.

In \cref{Sec:Gen}, we extend our results to the $RO(G)$-graded and Mackey functor-valued settings. 

\subsection{Conventions}

Throughout, $G$ denotes a finite group. If $M$ is an abelian group, then $\uM$ denotes the constant Mackey functor on $M$. 
In this work, ``manifold'' means ``smooth manifold.''

\subsection{Acknowledgments}

The authors thank Dan Dugger, Jeremy Hahn, Ben Knudsen, Irakli Patchkoria, Oscar Randal-Williams, Bridget Schreiner, Inna Zakharevich, and Mingcong Zeng for helpful discussions.
The first author was partially supported by NSF grant DMS-2204357.
The second author was partially supported by NSF grants DMS-2039316 amd DMS-2314082, as well as an AMS-Simons Travel Grant. The second author also thanks the Max Planck Institute in Bonn for providing a wonderful working environment and financial support during the beginning of this project.

\section{Equivariant homotopy theory}\label{Sec:Eqvt}

In this section, we discuss the ideas from equivariant homotopy theory which we will need in the sequel. \cref{SS:Bredon} contains a rapid recollection of Bredon homology with coefficients in a Mackey functor, which is the central equivariant homology theory we use throughout this work. \cref{SS:GFP} and \cref{SS:PHO} recall two standard constructions from equivariant stable homotopy theory: geometric fixed points and proper homotopy orbits. These auxiliary constructions appear in the isotropy separation sequence, which we recall in \cref{SS:ISS}. Our main technical result, \cref{Lemma:Bredon-to-PhiG}, applies the isotropy separation sequence to pass from Bredon homology to a simpler nonequivariant homology theory. This simpler theory is identified as a sum of singular homology theories in \cref{SS:GEM}. 

\subsection{Bredon homology}\label{SS:Bredon}

Bredon homology associates to a $G$-space a sequence of Mackey functors (see \cite[\S IX.4]{May96}). As in nonequivariant homotopy theory, Bredon homology with coefficients in a Mackey functor $\uM$ is representable in the equivariant stable homotopy category. By \cite[Theorem XIII.4.1]{May96}, there is an Eilenberg--MacLane $G$-spectrum $H\uM$ with the property that its homotopy Mackey functors satisfy $\upi_0(H\uM) \cong \uM$ and $\upi_i(H\uM)=0$ for $i>0$. If $X$ is a $G$-space, then we have
$$\uH_*^{(-)}(X;\uM) \cong \upi_*(\Sigma^\infty_G X \wedge H\uM),$$
where $\Sigma^\infty_G X$ is the \emph{suspension $G$-spectrum of $X$}, and $- \wedge - $ is the smash product of $G$-spectra. 

The Bredon homology of a $G$-CW complex can equivalently be defined using an equivariant analogue of the cellular chain complex \cite[\S I.4]{May96}. A \emph{coefficient system} is a contravariant functor $\co_G \to \Ab$, where $\co_G$ is the \emph{orbit category} of $G$. If $X$ is a $G$-CW complex, we may define a coefficient system by
$$\uC_n(X)(G/H) := \uH_n(X^n,X^{n-1}; \z)(G/H) := H_n((X^n)^H,(X^{n-1})^H; \z).$$
The connecting homomorphisms for the triples $((X^n)^H, (X^{n-1})^H, (X^{n-2})^H)$ define a map of coefficient systems
$$d: \uC_n(X) \to \uC_{n-1}(X).$$
If $\uM$ is a \emph{Mackey functor}, we may define cellular chains with coefficients in $\uM$ by
$$C_n^G(X;\uM) := \uC_n(X) \otimes \uM := \bigoplus_{G/H \in \co_G} \left( \uC_n(X)(G/H) \otimes \uM(G/H) \right)/\sim,$$
where $f^*c\otimes m \sim c \otimes f_*m$ for a map $f: G/H \to G/K \in \co_G$ and elements $c \in \uC_n(X)(G/H)$ and $m \in \uM(G/H)$. This is a chain complex of abelian groups, with differential given by $\partial = d \otimes 1$. By forgetting structure, we may analogously define 
$$C_n^H(X;\uM) := \uC_n(i^*_HX) \otimes i^*_H \uM$$
for each subgroup $H < G$; these assemble into a chain complex of coefficient systems whose homology
$$\uH_*^{(-)}(X;\uM) := H_*(C_n^{(-)}(X;\uM))$$
is the \emph{Bredon homology of $X$ with coefficients in $\uM$}. 

For much of the paper, we will be concerned with $\Z$-graded Bredon homology \emph{groups}. By this, we mean the $G/G$-levels $\uH_*^{G}(X;\uM)$ of the Bredon homology Mackey functors constructed here. We return to the level of generality of Bredon homology Mackey functors in \cref{Sec:Gen}. It is also possible to extend Bredon homology with coefficients in a Mackey functor from a $\Z$-graded collection of groups to a $RO(G)$-graded collection of groups, where $RO(G)$ denotes the representation ring of $G$ \cite[\S IX.5]{May96}. We review this extension in \cref{Sec:Gen} as well. 

\subsection{Geometric fixed points}\label{SS:GFP}

As outlined in the previous section, the Bredon homology groups of a $G$-space $X$ are the (categorical) fixed points of the $G$-spectrum $\sus^\infty_G X\sma H\uM$. The \emph{geometric fixed points} construction is related to categorical fixed points (in a way that will be made precise in \cref{SS:ISS}), but has properties that make it more amenable to computation. We refer the reader to \cite[Sec. 2.5]{HHR16} for proofs of the results in this section. 

\begin{definition} \label{Def:PhiG}
Let $G$ be a finite group and let $X$ be a $G$-spectrum. Let $\mathcal{P}$ be the family of proper subgroups of $G$, and $E\mathcal{P}$ be the unique $G$-homotopy type such that
$$ E\mathcal{P}^H \simeq \begin{cases} \emptyset & H=G, \\ * & H<G. \end{cases} $$
Consider the collapse map $E\mathcal P_+\to S^0$ and denote by $\widetilde{E\mathcal P}$ its cofiber. The \emph{geometric fixed points} of $X$ are given by
$$\Phi^GX = (\widetilde{E\mathcal{P}} \wedge X)^G. $$
\end{definition}

The key properties of the geometric fixed points functor are that it commutes with taking suspension spectra and is strong symmetric monoidal. These properties are the content of the following propositions. 

\begin{proposition}
For any $G$-space $X$, there is an equivalence of spectra
$$\Phi^G \Sigma^\infty_G X \simeq \Sigma^\infty(X^G).$$
\end{proposition}

\begin{proposition}
For any two $G$-spectra $X$ and $Y$, there is an equivalence of spectra
$$\Phi^G( X \wedge Y) \simeq \Phi^G X \wedge \Phi^G Y.$$
\end{proposition}

%Before stating the last key property, we need to recall an algebraic incarnation of geometric fixed points due to Blumberg--Gerhardt--Hill--Lawson \cite{BGHL19}. 

%\begin{definition}[{\cite[Defs. 5.4, 5.5]{BGHL19}}]
%Fix a finite group $G$ and let $N \subset G$ be a normal subgroup. Let $E\cf_N(\uA)$ denote the sub-Mackey functor of $\uA$ generated by $\uA(G/H)$ for all subgroups $H$ which do not contain $N$. Finally, let 
%$$\widetilde{E\cf}_N(\uA) := \uA/E\cf_N(\uA).$$
%If $\uM$ is a $G$-Mackey functor and $N$ is a normal subgroup of $G$, then let
%$$\widetilde{E\cf}_N\uM := \uM \boxtimes \widetilde{E\cf}_N\uA,$$
%where $- \boxtimes -$ denotes the box product of Mackey functors. 
%\end{definition}

%In particular, it follows from the definition that
%$$\widetilde{E\cf}_N(\uA)(G/H) = 
%\begin{cases}
%0 \quad & \text{ if } N \not\subset H, \\
%\uA((G/N)/(H/N)) \quad & \text{ if } N \subset H.
%\end{cases}
%$$

%In the case $\uM = \mfz$ and $N=G$, we can describe this construction more explicitly. 

%\begin{proposition}
%There is an isomorphism of abelian groups
%$$\widetilde{E\cf}_G\mfz(G/G) \cong
%\begin{cases}
%\f_p \quad & \text{ if $G$ is a $p$-group}, \\
%0 \quad & \text{ if $G$ is not a $p$-group}. 
%\end{cases}
%$$
%\end{proposition}

%\begin{proof}
%We wish to describe
%$$\widetilde{E\cf}_G\mfz(G/G) = (\mfz \boxtimes \widetilde{E\cf}_G(\uA))(G/G).$$
%According to \cite[Rmk. 2.13]{Zen17}, there is an isomorphism of abelian groups
%$$(\mfz \boxtimes \uM)(G/H) \cong \uM(G/H)/([H:K]x - \tr_K^H \res_K^H(x))$$
%for all Mackey functors $\uM$. Taking $H=G$ and $\uM=\widetilde{E\cf}_G(\uA)$, we have
%$$(\mfz \boxtimes \widetilde{E\cf}_G(\uA))(G/G) \cong \widetilde{E\cf}_G(\uA)(G/G)/([G:K]x - \tr_K^G\res_K^G(x)) \cong \z/([G:K]x - \tr_K^G\res_K^G(x)).$$
%Since $\widetilde{E\cf}_G(\uA)(G/G)$ is the quotient of $\uA$ obtained by killing transfers from proper subgroups, the relation above simplifies to just $[G:K]x=0$, so we have
%$$\widetilde{E\cf}_G\mfz(G/G) \cong \z/([G:K]x),$$
%where $K$ ranges over all proper subgroups of $G$. 

%If $G$ is a $p$-group, then we may take $K$ to be an index $p$ subgroup. Since all other subgroups have index a power of $p$, we obtain the first case. On the other hand, if $G$ is not a $p$-group, then we may take $K$ to be a $p$-Sylow subgroup and $K'$ to be a $q$-Sylow subgroup, where $p$ and $q$ are distinct primes. The quotient is then zero, which gives the second case. 
%\end{proof}

%\begin{proposition}[{\cite[Prop. 5.6]{BGHL19}}]
%For any Mackey functor $\uM$, we have
%$$\upi_0(\widetilde{E\cf}_N \wedge H\uM) \cong \widetilde{E\cf}_N\uM,$$
%where $\widetilde{E\cf}_N$ on the left denotes a cofibrant model of the universal space for the family of subgroups not containing $N$. 
%\end{proposition}

%Using that $\upi_0(\widetilde{E\cf}_G \wedge H\uM)(G/G) = \pi_0 \Phi^G H\uM$ (by definition), we obtain the following corollary. 

%\begin{corollary}\label{Cor:GFPZ}
%There is an isomorphism of abelian groups
%$$\pi_0\Phi^GH\mfz \cong
%\begin{cases}
%\f_p \quad & \text{ if } G \text{ is a $p$-group},\\
%0 \quad & \text{ if } G \text{ is not a $p$-group}. 
%\end{cases}$$
%\end{corollary}

%We apply \cref{Cor:GFPZ} to make the following computation. We thank Andrew Senger for suggesting the use of the Hopkins--Mahowald Theorem. 

%\begin{proposition}\label{Prop:GFPSplit}
%Let $G$ be a $p$-group and let $H\mfz$ denote the Eilenberg--MacLane $G$-spectrum for the constant Mackey functor $\mfz$. The geometric fixed points $\Phi^G H\mfz$ split as a wedge of nonnegative suspensions of $H\f_p$. 
%\end{proposition}

%\begin{proof}
%Since $G$ is a $p$-group, \cref{Cor:GFPZ} implies that $\pi_0\Phi^GH\mfz \cong \f_p$. Since $\mfz$ is a Green functor, $H\mfz$ is an $E_\infty$-ring spectrum, and since $\Phi^{C_p}(-)$ is strong symmetric monoidal, $\Phi^{C_p}H\mfz$ is an $E_\infty$-ring spectrum. Putting these facts together, we see that $\Phi^{C_p}H\mfz$ is an $E_2$-ring spectrum in which $p=0$. 

%The Hopkins--Mahowald Theorem (reinterpreted by Antol{\'i}n-Camarena--Barthel \cite[Thm. 5.1]{ACB19}) implies that $H\f_p$ is the initial $E_2$-ring spectrum in which $p=0$. Therefore there is a map of $E_2$-ring spectra 
%$$H\f_p \to \Phi^{C_p} H\mfz,$$
%and this equips $\Phi^{C_p} H\mfz$ with an $H\f_p$-module structure. Therefore $\Phi^{C_p}H\mfz$ splits as a wedge of suspensions of $H\mfp$ since $H\mfp$ is a field spectrum \cite[Sec. 1.3]{HS98}. Since geometric fixed points send connective $G$-spectra to connective spectra and $H\mfz$ is a connective $G$-spectrum, we see that $\Phi^G \mfz$ is a wedge of \emph{nonnegative} suspensions of $H\f_p$. 
%\end{proof}

%\begin{remark}
%When $G = (\z/p)^n$, the spectrum $\Phi^{G}H\mfz$ can be completely described using \cite[Thm. 3]{Kri20}. 
%\end{remark}

%\begin{proposition}\label{Prop:GFPVanish}
%If $G$ is not a $p$-group, then 
%$$\Phi^G H\mfz \simeq *.$$
%\end{proposition}

%\begin{proof}
%\cref{Cor:GFPZ} implies that $\pi_0 \Phi^G H\mfz = 0$; since $\pi_* \Phi^G H\mfz$ is a graded module over $\pi_0 \Phi^G H\mfz$, we have $\pi_* \Phi^G H\mfz = 0$. The result follows. 
%\end{proof}

%\begin{remark}
%\cref{Prop:GFPVanish} also follows from \cite[Prop. 11]{Kri20}, which shows that $\pi_*\Phi^G H\uA =0$ when $G$ is not a $p$-group, where $\uA$ is the Burnside Mackey functor. Since any Mackey functor is an $\uA$-module and $\Phi^G$ is symmetric monoidal, $\Phi^G H\mfz$ is a module over $\Phi^G H\uA = *$, and thus it is contractible. 
%\end{remark}


%\begin{proposition}\label{Prop:Green}
%Let $\uM$ be a Green functor. Then $\Phi^G H\uM$ splits as a wedge of nonnegative suspensions of Eilenberg--MacLane spectra. 
%\end{proposition}

%\begin{proof}
%Since $\uM$ is a Green functor, its value at $G/e$ is a commutative ring, and thus $i^*_e H\uM$ is an $H\z$-algebra. The unit map 
%$$H\z \to i^*_e H\uM$$
%is adjoint to a map of ring $G$-spectra
%$$i_* H\z \to H\uM$$
%via \cite[Prop. V.3.4]{MM02}. Applying geometric fixed points then yields a ring map (since $\Phi^G$ is strong symmetric monoidal)
%$$H\z \simeq \Phi^G i_* H\z \to \Phi^G H\uM,$$
%so $\Phi^G H\uM$ is an $H\z$-module. The result follows from the classical fact that any $H\z$-module is equivalent to a wedge of Eilenberg--MacLane spectra, plus geometric fixed points preserve connectivity. 
%\end{proof}

\subsection{Proper homotopy orbits}\label{SS:PHO}

In order to relate geometric fixed points to categorical fixed points, we will need the following auxiliary construction.

\begin{definition}
The \emph{proper homotopy orbits} of a $G$-spectrum $X$ are given by
$$X_{h\mathcal{P}} = (E\mathcal{P}_+ \wedge X)^G.$$
\end{definition}

% The homotopy groups of the proper homotopy orbits can be understood using a generalization of the homotopy orbit spectral sequence.

% \begin{proposition}[{\cite[Theorem 22.7]{GM95}; see also \cite[\S 3.1]{MNN19} and \cite{QS21}}]
% Let $X$ be a $G$-spectrum with $\pi_*^GX$ bounded below. There is a strongly convergent spectral sequence of abelian groups
% $$E^2_{p,q} = H_p^{G}(E\mathcal{P}; \pi_q^{(-)}X) \Rightarrow \pi_{p+q} X_{h\cp}.$$
% \end{proposition}

% \begin{remark}
% If $G = C_p$, then the family of proper subgroups is the family containing only the trivial subgroup, so the proper homotopy orbits are just the usual homotopy orbits, i.e.
% $$X_{h\cp} = X_{hC_p}.$$
% In this case, the proper homotopy orbit spectral sequence coincides with the usual homotopy orbit spectral sequence of \cite[\S 10]{GM95}. 

% More generally, the $E^2$-term of the proper homotopy orbit spectral sequence depends only on the restriction of the Mackey functor $\pi_q^{(-)}X$ to proper subgroups. 
% \end{remark}

%\EvaNote{I don't think this gets used?}
%\begin{lemma}
%Suppose $\pi_q^HX$ is finitely generated (as an abelian group) for each subgroup $H \leq G$. The Amitsur--Dress homology group $H_p^{\mathcal{P}}(G; \pi_q^G X)$ is finitely generated as an abelian group. 
%\end{lemma}

%\begin{proof}
%This follows from the description of Amitsur--Dress homology as the homology of a chain complex \cite[Eqn. 20.4]{GM95} along with our finiteness assumption on $\pi_q^HX$. 
%\end{proof}


\subsection{Isotropy separation sequence}\label{SS:ISS}

In this section, we prove \cref{Lemma:Bredon-to-PhiG}, a central tool in our analysis of Bredon homology. It is based on the following cofiber sequence, which allows us to understand categorical fixed points, provided we can understand geometric fixed points and proper homotopy orbits.

\begin{lemma}[Isotropy Separation Sequence {\cite[Section 2.5.2]{HHR16}}]
The collapse map of \cref{Def:PhiG} induces a cofiber sequence of spectra
$$ X_{h\mathcal P} \to X^G \to \Phi^G X. $$
\end{lemma}

% \begin{lemma}[Five Lemma mod Serre]
% Suppose we have a commutative diagram
% \[
% \begin{tikzcd}
% A \arrow{r}{f} \arrow{d}{\ell} & B \arrow{r}{g} \arrow{d}{m} & C \arrow{r}{h} \arrow{d}{n} & D \arrow{r}{j} \arrow{d}{p} & E \arrow{d}{q} \\
% A' \arrow{r}{r} & B' \arrow{r}{s} & C' \arrow{r}{t} & D' \arrow{r}{u} & E'
% \end{tikzcd}
% \]
% in an abelian category $\mathcal{A}$. Let $\mathcal{C} \subseteq \mathcal{A}$ be a Serre subcategory. If the rows are exact, $m$ and $p$ are isomorphisms mod $\mathcal{C}$, $\ell$ is an epimorphism mod $\mathcal{C}$, and $q$ is a monomorphism mod $\mathcal{C}$, then $n$ is an isomorphism mod $\mathcal{C}$. 
% \end{lemma}

% \begin{proof}
% As $\mathcal{C}$ is a Serre subcategory, there exists an abelian category $\mathcal{A}/\mathcal{C}$ and an exact functor $F: \mathcal{A} \to \mathcal{A}/\mathcal{C}$ which is essentially surjective with kernel $\mathcal{C}$ \cite[Lem. 12.10.6]{Stacks}. Applying $F$ yields a commutative diagram
% \[
% \begin{tikzcd}
% FA \arrow{r}{Ff} \arrow{d}{F\ell} & FB \arrow{r}{Fg} \arrow{d}{Fm} & FC \arrow{r}{Fh} \arrow{d}{Fn} & FD \arrow{r}{Fj} \arrow{d}{Fp} & FE \arrow{d}{Fq} \\
% FA' \arrow{r}{Fr} & FB' \arrow{r}{Fs} & FC' \arrow{r}{Ft} & FD' \arrow{r}{Fu} & FE'
% \end{tikzcd}
% \]
% in the abelian category $\mathcal{C}/\mathcal{A}$ in which the rows are exact, $Fm$ and $Fp$ are isomorphisms, $F\ell$ is an epimorphism, and $Fq$ is a monomorphism. The classical Five Lemma implies that $Fn$ is an isomorphism, so $F\ker(n)=\ker(Fn)=0$ and $F\coker(n)=\coker(Fn)=0$, i.e., $\ker(n)$ and $\coker(n)$ lie in $\mathcal{C}$, i.e., $n$ is an isomorphism mod $\mathcal{C}$. 
% \end{proof}

%\begin{proposition}\label{Prop:Bredon-to-PhiG}
%Let $G$ be a finite group and let $X, Y \in \Sp^G$. Suppose that there is a map of $G$-spectra $f : X \to Y$ which induces an isomorphism 
%$$H_i(X^e; \z) \xrightarrow{\cong} H_i(Y^e; \z)$$
%for all $i \leq d$. Then if
%$$\coker( \pi_d (\Phi^G( H\mfz \wedge X)) \xrightarrow{f_*} \pi_d(\Phi^G(H\mfz \wedge Y)))$$
%is a finitely generated abelian group,
%$$\coker( H_d^G(X; \mfz) \xrightarrow{f_*} H_d^G(Y; \mfz) )$$
%is a finitely generated abelian group. 
%\end{proposition}

%\begin{lemma}\label{Lemma:Bredon-to-PhiG}
%Let $G$ be a finite group and let $X\in \Sp^G$. Let $\sigma_1,\dots,\sigma_k$ be equivariant self-maps $X\to X$ such that 
%\begin{equation}\label{Eqn:ADCondition}
%\coker\left( (\sigma_k)_*: H_i^G(E\mathcal{P}; H_j^{(-)}(X; \mfz)) \to H_i^G(E\mathcal{P}; H_j^{(-)}(X; \mfz)) \right)
%\end{equation}
%is a finitely generated $\Z[(\sigma_1)_*,\ldots,(\sigma_{k-1})_*]$-module for all $i,j \leq d$. If
%$$\coker\left( \pi_i (\Phi^G( H\mfz \wedge X)) \xrightarrow{(\sigma_k)_*} \pi_i(\Phi^G(H\mfz \wedge X)) \right)$$
%is a finitely generated $\Z[(\sigma_1)_*,\dots,(\sigma_{k-1})_*]$-module for all $i \leq d$, then
%$$\coker\left( H_d^G(X; \mfz) \xrightarrow{(\sigma_k)_*} H_d^G(X; \mfz) \right)$$
%is a finitely generated $\Z[(\sigma_1)_*,\dots,(\sigma_{k-1})_*]$-module.
%\end{lemma}

\begin{lemma}\label{Lemma:Bredon-to-PhiG}
Let $G$ be a finite group, let $d\geq 0$, and let $X$ be a $G$-spectrum. Let
$\mathcal{C}$ be a Serre class. Assume $\pi_q \Phi^K(H\underline{\Z}\sm X)\in
\mathcal{C}$ for all $K\leq G$, $q\leq d+1$. Then $H_q^K(X;\underline{\Z})\in
\mathcal{C}$ for all $K\leq G$, $q\leq d$.
\end{lemma}

The proof relies on the following homological fact, the proof of which we learned from Piotr Pstr{\k{a}}gowski. 

\begin{lemma}[Four Lemma mod Serre]\label{Lem:FourModC}
Suppose we have a commutative diagram
\[
\begin{tikzcd}
A \arrow{r}{f} \arrow{d}{\ell} & B \arrow{r}{g} \arrow{d}{m} & C \arrow{r}{h} \arrow{d}{n} & D \arrow{d}{p} \\
A' \arrow{r}{r} & B' \arrow{r}{s} & C' \arrow{r}{t} & D'
\end{tikzcd}
\]
in an abelian category $\mathcal{A}$. Let $\mathcal{C} \subseteq \mathcal{A}$ be a Serre subcategory. Suppose the rows are exact, $\ell$ is an epimorphism mod $\mathcal{C}$ and $p$ is a monomorphism mod $\mathcal{C}$. If $n$ is an epimorphism mod $\mathcal{C}$ then so is $m$. If $m$ is a monomorphism mod $\mathcal{C}$ then so is $n$.
\end{lemma}

\begin{proof}
We just prove the first statement, as the second is analogous.
As $\mathcal{C}$ is a Serre subcategory, there exists an abelian category $\mathcal{A}/\mathcal{C}$ and an exact functor $F: \mathcal{A} \to \mathcal{A}/\mathcal{C}$ which is essentially surjective with kernel $\mathcal{C}$ \cite[Lemma 02MS]{Stacks}. Applying $F$ yields a commutative diagram
\[
\begin{tikzcd}
FA \arrow{r}{Ff} \arrow{d}{F\ell} & FB \arrow{r}{Fg} \arrow{d}{Fm} & FC \arrow{r}{Fh} \arrow{d}{Fn} & FD \arrow{d}{Fp}\\
FA' \arrow{r}{Fr} & FB' \arrow{r}{Fs} & FC' \arrow{r}{Ft} & FD' 
\end{tikzcd}
\]
in the abelian category $\mathcal{C}/\mathcal{A}$ in which the rows are exact, $F\ell$ and $Fn$ are epimorphisms, and $Fp$ is a monomorphism. The classical Four Lemma implies that $Fm$ is an epimorphism, so $F\coker(m)=\coker(Fm)=0$, i.e., $\coker(m)$ lies in $\mathcal{C}$, i.e., $m$ is an epimorphism mod $\mathcal{C}$. 
\end{proof}

\begin{lemma} \label{lem:sseq-iso-mod-serre}
Let $\mathcal{C}\subseteq \mathcal{A}$ be a Serre category.
Let $\sigma:E'_r \to E''_r$ be a map of spectral sequences that is an isomorphism modulo $\mathcal{C}$ on the $E_r$ page for some $r$. Then $E'_{r+1} \too{\sigma} E''_{r+1}$ is an isomorphism mod $\mathcal{C}$.
\end{lemma}

\begin{proof}
This is a straightforward consequence of Lemma \ref{Lem:FourModC}. First apply Lemma \ref{Lem:FourModC} to the diagram
$$ \xymatrix{
0\ar[r]\ar[d] & Z'_r\ar[r]\ar[d] & E'_r\ar[r]^-{d_r}\ar[d]^\sigma & E'_r\ar[d]^\sigma
\\0\ar[r] & Z''_r\ar[r] & E''_r\ar[r]^-{d_r} & E''_r
}$$
(where $Z'_r$ denotes the submodule of cycles) to show that $Z'_r\to Z''_r$ is an isomorphism mod $\mathcal{C}$. Similar arguments involving the exact sequences $0\to Z'_r \to E'_r \to B'_r \to 0$ and $0\to B'_r \to Z'_r \to Z'_r/B'_r= E'_{r+1}\to 0$ show that the induced maps $Z'_r\to Z''_r$ and $E'_{r+1}\to E''_{r+1}$, respectively, are isomorphisms mod $\mathcal{C}$.
\end{proof}

\begin{proof}[Proof of \cref{Lemma:Bredon-to-PhiG}]
By induction up the subgroup lattice. For the base case, 
$$ H^e_q(X; \underline{\Z}) = \pi_q((H\underline{\Z}\sm X)^e) =
\pi_q(\Phi^e(H\underline{\Z}\sm X)) \in \mathcal{C}. $$
Now let $\{ e \} < K\leq G$ and let $\mathcal{P}$ denote the family of proper
subgroups of $K$.
The isotropy separation sequence
$$ (H\underline{\Z}\sm X)_{h\mathcal{P}}\to
(H\underline{\Z}\sm X)^K\to \Phi^K(H\underline{\Z}\sm X)$$
gives rise to a long exact sequence
\begin{gather*}
\cdots \to \pi_{q+1}(\Phi^K(H\underline{\Z}\sm X))\to \pi_q((H\underline{\Z}\sm X)_{h\mathcal{P}}) \to H^K_q(X; \underline{\Z}) 
\\\to \pi_q(\Phi^K(H\underline{\Z}\sm X)) \to \pi_{q-1}((H\underline{\Z}\sm X)_{h\mathcal{P}}) \to \cdots. 
\end{gather*}
To show that the middle term is in $\mathcal{C}$, we claim it suffices to show that all
the other terms are in $\mathcal{C}$: apply Lemma \ref{Lem:FourModC}
to the map from the above exact sequence to the zero exact sequence.

The assumption says that the first and fourth terms are in $\mathcal{C}$, so we
focus on showing the second term is in $\mathcal{C}$, which would also show the fifth
term is in $\mathcal{C}$ because $q\leq d$ is arbitrary.
There is a cell decomposition of $E\mathcal{P}_+$ by a finite collection of
cells
$$ \{ {K/H_i}_+ \sm S^{n_i} \}_{i\in \mathcal{I}} $$
with every $H_i < K$ (cf. \cite[Section 2.5.2]{HHR16}). Filtering $E\mathcal{P}_+\sm X$ by the cells of
$E\mathcal{P}_+$ gives rise to a spectral sequence
$$ E_1^{i,p-n_i} \cong H_p^K({K/H_i}_+\sm S^{n_i}\sm X)
\Rightarrow H^K_{p-n_i}(E\mathcal{P}_+\sm X). $$
We have
\begin{align*}
H^K_p({K/H}_+\sm S^n\sm X)  & \cong [S, H\underline{\Z}\sm K/H\sm X]^K_{p-n}
\cong [{K/H}_+, H\underline{\Z}\sm X]^K_{p-n}
\\ & \cong [S, H\underline{\Z}\sm X]^H_{p-n} \cong H^H_{p-n}(X;\underline{\Z})
\end{align*}
using self-duality of $K/H_+$. In particular, since $H_i<K$, 
we have $E_1^{i,q}\in \mathcal{C}$ by the induction hypothesis. By Lemma
\ref{lem:sseq-iso-mod-serre}, $H^K_q(E\mathcal{P}_+\sm X)$ is also in
$\mathcal{C}$.
\end{proof}


%\begin{proof}[Proof of \cref{Lemma:Bredon-to-PhiG}]
%Naturality of proper homotopy orbits, categorical fixed points, and geometric fixed points implies that $\sigma_k$ induces a map of cofiber sequences of spectra 
%\[
%\begin{tikzcd}
%(H\mfz \wedge X)_{h\cp} \arrow{r} \arrow{d} & (H \mfz \wedge X)^G \arrow{r} \arrow{d} & \Phi^G(H\mfz \wedge X) \arrow{d} \\
%(H\mfz \wedge X)_{h\cp} \arrow{r} & (H\mfz \wedge X)^G \arrow{r} & \Phi^G(H\mfz \wedge X).
%\end{tikzcd}
%\]
%Applying homotopy groups gives a map of long exact sequences
%\[
%\begin{tikzcd}
%\pi_d(H\mfz \wedge X)_{h\cp} \arrow{r} \arrow{d}{\alpha_d} & H_d^G(X;\mfz) \arrow{r} \arrow{d} & \pi_d\Phi^G(H\mfz \wedge X)  \arrow{d} \arrow{r} & \pi_{d-1}(H\mfz \wedge X)_{h\cp} \arrow{d}{\alpha_{d-1}} \\
%\pi_d(H\mfz \wedge X)_{h\cp} \arrow{r} & H_d^G(X;\mfz) \arrow{r} & \pi_d\Phi^G(H\mfz \wedge X)  \arrow{r} & \pi_{d-1}(H\mfz \wedge X)_{h\cp}.
%\end{tikzcd}
%\]
%The result will follow from \cref{Lem:FourModC} applied to the Serre class of finitely generated $\Z[(\sigma_1)_*,\dots,(\sigma_{k-1})_*]$-modules after we show that the maps $\alpha_d$ and $\alpha_{d-1}$ are isomorphisms mod $\mathcal{C}$.
%
%To see this, we apply the proper homotopy orbit spectral sequence. Naturality yields a map of spectral sequences
%\[
%\begin{tikzcd}
%E^2_{p,q} = H_p^{G}(E\mathcal{P}; H_q^{(-)}(X;\mfz)) \arrow[Rightarrow, r] \arrow{d} & \pi_{p+q}(H\mfz \wedge X)_{h\cp} \arrow{d}{\alpha_{p+q}} \\
%E^2_{p,q} = H_p^{G}(E\mathcal{P}; H_q^{(-)}(X;\mfz)) \arrow[Rightarrow, r] & \pi_{p+q}(H\mfz \wedge X)_{h\cp} .
%\end{tikzcd}
%\]
%Restricting to total degree $d$, we have
%\[
%\begin{tikzcd}
%\bigoplus_{p+q=d} H_p^{G}(E\mathcal{P}; H_q^{(-)}(X;\mfz)) \arrow[Rightarrow, r] \arrow{d} & \pi_d^G(H\mfz \wedge X)_{h\cp} \arrow{d}{\alpha_d} \\
%\bigoplus_{p+q=d} H_p^{G}(E\mathcal{P}; H_q^{(-)}(X;\mfz)) \arrow[Rightarrow, r] & \pi_d^G(H\mfz \wedge X)_{h\cp}.
%\end{tikzcd}
%\]
%
%Since we have assumed that $\sigma_k$ induces an isomorphism on integral homology in degrees $0 \leq i \leq d$ mod $\mathcal{C}$, the map on the left is an isomorphism mod $\mathcal{C}$ on $E^2$-terms. Since $\sigma_k$ induces a map of spectral sequences, Lemma \ref{lem:sseq-iso-mod-serre} implies that $\alpha_d$ is also an isomorphism mod $\mathcal{C}$. The same argument applies if we replace $d$ by $d-1$, so the result follows.
%\end{proof}
%
%\begin{remark}\label{Rmk:ADWeak}
%Assuming the other hypotheses of \cref{Lemma:Bredon-to-PhiG}, the condition \eqref{Eqn:ADCondition} follows from the weaker assumption that
%\begin{equation}\label{Eqn:ADWeak}
%\coker \left( (\sigma_k^e)_* : H_i(X;\z) \to H_i(X;\z) \right)
%\end{equation}
%is a finitely generated $\Z[(\sigma_1^e)_*,\ldots,(\sigma_{k-1}^e)_*]$-module for all $i \leq d$. Note that this weaker assumption holds when $\sigma_k^e : X \to X$ is the classical stabilization map and $X = C(M) = \bigsqcup_{j \geq 0} C_j(M)$ is the disjoint union of unordered configuration spaces of a stabilizable manifold. 
%
%The stronger assumption \eqref{Eqn:ADCondition} follows from the weaker assumption \eqref{Eqn:ADWeak} by induction up the subgroup lattice of $G$. Let $(*_H)$ be the statement
%\begin{equation}\label{Eqn:ADH}
%\tag{$\ast_H$}\atop{\coker\left( (\sigma_k)_* :
%H_i^H(E\mathcal{P}; H_j^{(-)}(X;\mfz)) \to H_i^H(E\mathcal{P};
%H_j^{(-)}(X;\mfz)) \right)
%\\\text{is a finitely generated
%$\z[(\sigma_1)_*,\ldots,(\sigma_{k-1})_*]$-module for all $i,j\leq d$.}
%}
%\end{equation}
%Then our weaker assumption \eqref{Eqn:ADWeak} is $(*_e)$ and we wish to prove the stronger statement \eqref{Eqn:ADCondition}, the case $(*_G)$.
%
%We proceed by induction up the subgroup lattice. The base case is $H=e$, which is \eqref{Eqn:ADWeak}. Suppose now that $(*_K)$ holds for all proper subgroups $K<H$ and that we wish to prove $(*_H)$. The groups $H_i^H(E\mathcal{P}; H_j^{(-)}(X;\mfz))$ depend only on the restriction of the $H$-Mackey functor $H_j^{(-)}(X;\mfz)$ to proper subgroups of $H$, so it suffices to show that the cokernel of the map on coefficients
%$$\coker\left( (\sigma_k)_*: H_j^{K}(X;\mfz) \to H_j^{K}(X;\mfz) \right)$$
%is a finitely generated $\z[(\sigma_1)_*,\ldots,(\sigma_{k-1})_*]$-module for all proper subgroups $K < H$ and $j \leq d$. But since we have assumed $(*_K)$ holds for all proper subgroups $K < H$, this follows immediately by applying \cref{Lemma:Bredon-to-PhiG} for each $K$. 
%
%%The base case $G=e$ is \eqref{Eqn:ADWeak}. Let $e < H \leq G$ be a nontrivial subgroup with no intermediate subgroup between $e$ and $H$. Feeding the weaker assumption \eqref{Eqn:ADWeak} into \cref{Lemma:Bredon-to-PhiG} implies that 
%%$$\coker\left( H_i^{H}(X; \mfz) \to H_i^{H}(X; \mfz) \right)$$
%%is a finitely generated $\z[(\sigma_1)_*,\ldots,(\sigma_{k-1})_*]$-module for all $i \leq d$, which implies that
%%$$\coker \left( (\sigma_k)_*: H_i^H(E\mathcal{P}; H_t^{(-)}(X;\mfz)) \to H_i^H(E\mathcal{P}; H_t^{(-)}(X;\mfz)) \right)$$
%%is a finitely generated $\z[(\sigma_1)_*,\ldots,(\sigma_{k-1})_*]$-module for all $i \leq d$. Repeating this argument with $H$ in place of $e$ moves us further up the subgroup lattice, and iterating this process eventually yields \eqref{Eqn:ADCondition}. 
%
%%\jd{We could make this more precise by defining condition $*K*$ for each $K \leq G$ to be the analogue of the stronger assumption for the restriction to $K$, and then show that if $H < K$ is an adjacent proper subgroup, then $*H*$ implies $*K*$. But if the previous sketch is clear, I'd prefer not to get that technical.}
%\end{remark}


%\begin{proposition}[$G$-Whitehead mod Serre]
%Let $\cf$ be a family of subgroups of $G$ and let $\cc$ be the Serre subcategory of finitely generated abelian groups. Let $f: X \to Y$ be a map of connective $G$-spectra with $\coker(f^H) \in \cc$ for all $H \in \cf$. 
%Then the map
%$$f \wedge E\cf_{+} : X \wedge E\cf_+ \to Y \wedge E\cf_+$$
%is an equivalence of $G$-spectra mod $\cc$. 
%\end{proposition}

%\begin{proof}
%By definition, we must show that $f^H : (X \wedge E\cf_+)^H \to (Y \wedge E\cf_+)^H$ is an equivalence mod $\cc$ for all subgroups $H \leq G$. The isotropy separation sequence gives rise to a map of long exact sequences
%\[
%\begin{tikzcd}
%\cdots \arrow{r} \arrow{d} & \pi_{d+1}(X \wedge \widetilde{E\cf})^H \arrow{r} \arrow{d}{f \wedge \widetilde{E\cf}} & \pi_d(X \wedge E\cf_+)^H \arrow{r} \arrow{d}{f \wedge E\cf_+} & \pi_dX^H \arrow{r} \arrow{d}{f} & \cdots \\
%\cdots \arrow{r} & \pi_{d+1}(Y \wedge \widetilde{E\cf})^H \arrow{r} & \pi_d(Y \wedge E\cf_+)^H \arrow{r} & \pi_dY^H \arrow{r} & \cdots .
%\end{tikzcd}
%\]

%If $H \in \cf$, then the equivalences of $H$-spaces
%$$\widetilde{E\cf} \simeq *, \quad E\cf_+ \simeq S^0$$
%imply that $f \wedge E\cf_+$ is an equivalence mod $\cc$ if and only if $f$ is an equivalence mod $\cc$, and this holds by assumption. 

%If $H \notin \cf$, we proceed by induction up the subgroup lattice. Our induction hypothesis is that $\coker((f \wedge E\cf_+)^K) \in \cc$ for all proper subgroups $K < H$, and we wish to show that $\coker((f \wedge E\cf_+)^H) \in \cc$. Observe that
%$$(E\cf_+)^H \simeq *.$$
%Consider the hyper-Tor spectral sequence of \cite{LM06},
%$$E^2_{s,t} = \underline{\Tor}_{s,t}^{\pi_*^{(-)}(S^0)}(\pi_*^{(-)}(C), \pi_*^{(-)}(E\cf_+)) \Rightarrow \pi_{s+t}^{(-)}(C \wedge E\cf_+),$$
%where $C = \coker(f: X \to Y)$. 

%Since $C$ and $E\cf_+$ are connective, the induction hypothesis implies that (after resolving extensions) the portion of the $E^\infty$-term converging to $\pi_n^{K}(C \wedge E\cf_+)$ is in $\cc$ for each proper subgroup $K \leq H$. The only other part of the $E^2$-term which contributes to $\pi_n^H(C \wedge E\cf_+)$ is zero since $\pi_*^H(E\cf_+)$ is a trivial module. Therefore $\pi_n^H(C \wedge E\cf_+) \in \cc$. 
%\end{proof}

\subsection{Geometric fixed points of Eilenberg--MacLane $G$-spectra}\label{SS:GEM}

In light of the previous section, one piece of understanding the Bredon homology $(H\uM\sma \sus^\infty_GX)^G$ of a $G$-space $X$ is understanding a geometric fixed points term of the form
$$ \Phi^G(H\uM\sma \sus^\infty_GX) \simeq \Phi^G H\uM \sma \Phi^G \sus^\infty_GX \simeq \Phi^G H\uM \sma \sus^\infty X^G, $$ 
which is the $\Phi^G H\uM$-homology of the space $X^G$. In many cases of interest (see Proposition \ref{Prop:Green}), this splits into (nonequivariant) Eilenberg-Maclane spaces, allowing us to reduce to studying ordinary homology of fixed point spaces.

Recall that the category of Mackey functors is symmetric monoidal with unit the Burnside Mackey functor $\uA$ and monoidal product the \emph{box product}. A Green functor is a monoid in the category of Mackey functors \cite{lewis-green}.

\begin{proposition}\label{Prop:Green}
Let $\uM$ be a Green functor. For all $H \leq G$, the spectrum $\Phi^H H\uM$ is a connective generalized Eilenberg--MacLane spectrum, i.e., it splits as a wedge of nonnegative suspensions of Eilenberg--MacLane spectra. Moreover, if $\uM(G/K)$ is finitely generated as an abelian group for all $K \leq H$, then $\pi_k \Phi^H H\uM$ is a finitely generated abelian group for each $k \in \z$. 
\end{proposition}

\begin{proof}
The lax monoidal transformation $(-)^H \to \Phi^H(-)$ induces a map of commutative ring spectra (cf. \cite[\S B.10.5]{HHR16})
$$(H\uM)^H \to \Phi^H H\uM.$$
Since $H\uM$ is Eilenberg--MacLane, $(H\uM)^H \simeq H\uM(G/H)$, and since $\uM$ is a Green functor, $\uM(G/H)$ is a commutative ring (see, e.g., \cite[\S 2.2.5]{Shu10}). Therefore we have a map of commutative ring spectra
$$H\z \to H\uM(G/H) \simeq (H\uM)^H \to \Phi^H H\uM$$
which determines an $H\z$-module structure on $\Phi^H H\uM$. The claim that $\Phi^H H\uM$ is generalized Eilenberg--MacLane then follows from the classical fact (usually attributed to Adams) that any $H\z$-module is generalized Eilenberg--MacLane. 

That $\Phi^H H\uM$ is connective follows from the isotropy separation sequence: we have
$$\Phi^H H\uM \simeq \cofib((H\uM)_{h\mathcal{P}} \to (H\uM)^{H})$$
and both $(H\uM)_{h\mathcal{P}}$ and $(H\uM)^H$ are connective, so $\Phi^H H\uM$ is connective. 

Finite generation also follows from the isotropy separation sequence: the finite generation hypothesis on $\uM(G/K)$, along with the fact that $E\mathcal{P}$ admits an $H$-CW structure with finitely cells in each dimension, implies that each homotopy group of $(H\uM)_{h \mathcal{P}}$ and $(H\uM)^H \simeq H(\uM(G/H))$ is finitely generated, so the same is true for $\Phi^H H\uM$. 
\end{proof}

\begin{example}
The Burnside Mackey functor $\uA$ is the Mackey functor given by 
$$ \uA(G/K)=A(K), $$
where $A(K)$ denotes the Grothendieck ring of finite $K$-sets. The restrictions and transfers in $\uA$ are induced by restriction and induction of $K$-sets, respectively. 

The constant Mackey functor $\mfz$ has $\mfz(G/K)=\Z$ in every level. The restrictions in $\mfz$ are the identity map $\Z\to\Z$ and the transfer $\mfz(G/K)\to\mfz(G/L)$ is given by multiplication by the index $[L:K]$.

Both $\uA$ and $\mfz$ are Green functors, so $\Phi^H H\uA$ and $\Phi^H H\mfz$ are both connective generalized Eilenberg--MacLane spectra. Moreover, $\uA(G/K) = A(K)$ and $\mfz(G/K) = \z$ are finitely generated abelian groups for all $K \leq H$, so $\pi_k \Phi^H H\uA$ and $\pi_k \Phi^H H\mfz$ are both finitely generated abelian groups for all $k\geq 0$. 
\end{example}

\begin{remark}
It is possible to describe $\Phi^G H\uM$ explicitly in many cases of interest using previous computations of $\pi_*\Phi^G H\uM$ (e.g., \cite{HHR16, Kri20}), but this will not be necessary for our applications. 
\end{remark}


%The following proposition implies that 

%\begin{proposition}
%Let $\uR$ be a Green functor with $\uR(G/H)$ finitely generated for all $H \leq G$. Then
%$$\pi_n \Phi^G H\uR$$
%is a finitely generated group for all $n \in \z$. 
%\end{proposition}

%\begin{proof}
%The isotropy separation sequence
%$$(H\uR)_{h\mathcal{P}} \to (H\uR)^G \to \Phi^G H\uR$$
%gives rise to a long exact sequence in homotopy
%$$\cdots \to \pi_n(H\uR)_{h\mathcal{P}} \to \pi_n (H\uR)^G \to \pi_n \Phi^G H\uR \to \pi_{n-1}(H\uR)_{h\mathcal{P}} \to \cdots.$$
%We have 
%$$\pi_n (H\uR)^G \cong \begin{cases}
%\uR(G/G) \quad & \text{ if } n=0, \\
%0 \quad & \text{ otherwise,}
%\end{cases}$$
%and $\uR(G/G)$ is finitely generated by assuption. Therefore it suffices to show that $\pi_n(H\uR)_{h\mathcal{P}}$ is finitely generated for all $n$. This follows from the proper homotopy orbit spectral sequence: the associated graded of $\pi_n(H\uR)_{h\mathcal{P}}$ is a subquotient of 
%$$\bigoplus_{s+t=n} H_s^G(E\mathcal{P};\pi^{(-)}_t(H\uR)^G) \cong \bigoplus_{0 \leq s \leq n} H_s^G(E\mathcal{P}; \uR),$$
%which is a sum of finitely generated groups (cf. the proof of \cref{Prop:Green}). 
%\end{proof}

We conclude by recording the following fact which can be used to reduce to $p$-Sylow subgroups for coefficients in a constant Mackey functor. We will not need it later, but it can be used to simplify some arguments in the sequel if one works with constant Mackey functor coefficients. 

\begin{proposition}\label{Prop:GFPVanish}
Let $\uR$ be a constant Green functor and let $\uM$ be an $\uR$-algebra. If $G$ is not a $p$-group, then
$$\Phi^G H\uM \simeq *.$$
\end{proposition}


\begin{proof}
The $G$-spectrum $H\uM$ is an $H\uR$-algebra, and since geometric fixed points are strong symmetric monoidal, the spectrum $\Phi^G H\uM$ is a $\Phi^G H\uR$-algebra. But \cite[Prop. 11]{Kri20} implies that when $G$ is not a $p$-group, $\Phi^G H\uR \simeq *$, from which we conclude $\Phi^G H\uM \simeq *$. 
\end{proof}

\begin{example}
The constant Green functor $\mfz$ is an algebra over itself, so $\Phi^G H\mfz \simeq *$ if $G$ is not a $p$-group. On the other hand, one can show that $\pi_0 \Phi^G H\uA \cong \z$ for any finite group $G$, so $\Phi^G H\uA \not\simeq *$ for any $G$. 
\end{example}

\begin{remark}
For simplicity, we will work with $\mfz$ coefficients in the sequel. If $\uM$ is a constant Mackey functor which is levelwise finitely generated, then it is straightforward to deduce $\uM$ coefficient analogues of all of the $\mfz$ coefficient statements below using the presentation of $\uM$ as a $\mfz$-module (cf. deducing homology isomorphisms with $M$ coefficients from homology isomorphisms with $\z$ coefficients). 

For certain groups, e.g., any cyclic $p$-group $C_{p^n}$, it is also straightforward to pass from $\mfz$ coefficients to $\uA$ coefficients using induction up the subgroup lattice, the short exact sequence of Mackey functors
$$0 \to \uI \to \uA \to \mfz \to 0,$$
and the fact that $H_*^{C_{p^n}}(X;\uI) \cong H_*^{C_{p^{n-1}}}(X^{C_p};\uI')$, where $\uI'$ is the $C_{p^{n-1}}$ Mackey functor with $\uI'(C_{p^{n-1}}/C_{p^m}) = \uI(C_{p^n}/C_{p^{m+1}})$. Since every Mackey functor is an $\uA$-module, one can deduce analogues of our results with coefficients in any Mackey functor which is finitely generated for cyclic $p$-groups. 

We leave the extension to other coefficients and more complicated groups to the interested reader. 
\end{remark}


\section{Configuration spaces and stabilization maps}\label{Sec:Conf}

As in the previous section, $G$ is any finite group. In this section, we discuss configuration spaces and stabilization maps. We recall the nonequivariant theory in \cref{SS:Conf} and discuss the equivariant analogues in \cref{SS:CGS} and \cref{SS:GConf}. In \cref{SS:CGS}, we define the more nuanced notion of $S$-configuration spaces (\cref{Def:ConfG} and \cref{Def:CGS}) and their stabilization maps. Our main results are \cref{Prop:FixptConfDecomp}, which expresses the $G$-fixed points of configuration spaces in terms of $S$-configurations, and \cref{Cor:CGSDecomp} which expresses each $S$-configuration space in terms of products of ordinary configuration spaces. Combining these results yields \cref{MT:CtoCSG}.

\subsection{Nonequivariant configuration spaces and stabilization maps}\label{SS:Conf}


\begin{definition}\label{Def:Conf}
Let $M$ be a space.
The \emph{configuration space of $n$ ordered points in $M$} is
$$ \Conf_n(M):=\Emb(\{1,\ldots,n\},M)= \{ (x_1,\dots,x_n)\in M^{\x n} \st x_i \neq x_j \text{ for all
}i\neq j \} $$
topologized as a subspace of $M^{\x n}$.

The \emph{configuration space of $n$ unordered points in $M$} is the quotient space
$$C_n(M) := \Conf_n(M)/\Sigma_n,$$
where the symmetric group $\Sigma_n$ acts by permuting points. 
\end{definition}

%\begin{definition}
%Let $M$ and $X$ be spaces.
%\EvaNote{Can we omit $X$? Confused about this.}
%The \emph{configuration space of $n$ unordered points in $M$ with labels in $X$} is defined by
%$$C_n(M,X) := (\Emb(n,M) \times X^n)/\Sigma_n.$$
%\end{definition}

When $M$ is an open, connected manifold with $\dim M\geq 2$, then $M$ is the interior of a manifold with boundary, and there there are well-defined \emph{stabilization maps} $C_n(M) \to C_{n+1}(M)$ (cf. \cite[Secs. 2.1-2.2]{Pal18}). Since we will describe the equivariant analogue of these maps in some detail in \cref{SS:GConf}, we review the non-equivariant construction.

\begin{definition}
We say that an $n$-dimensional manifold $M$ is \emph{stabilizable} if $M$ is homeomorphic to the interior of a $n$-dimensional manifold with nonempty boundary. 
\end{definition}

\begin{construction}
\label{Constr:nonequivariant-stabilization}
When $M$ is stabilizable, we can define a stabilization map on unordered configurations as follows. Let $M$ be stabilizable with $M \cong W^\circ$, where $W$ is an $n$-dimensional manifold with nonempty boundary. Let $p \in \partial W$. By the definition of manifold with boundary, 
%\EvaNoteHl{collar neighborhood theorem,}{Do we need a theorem or is this just the definition manifold with boundary? (boundary version of Euclidean neighborhood)}
there exists an open neighborhood $U \subset W$ containing $p$ together with a diffeomorphism 
$$\phi: U \xrightarrow{\cong} \r^n_+ = \{ x = (x_1,\ldots,x_n) \in \r^n : x_1 \geq 0\}$$ 
sending $p$ to $0$. Let $b: \r^n_+ \to \r^n_+$ be a smooth bump function which  sends $0$ to $(1,0,\ldots,0)$. Let $e: W \to W$ be the self-embedding defined by
$$e(x) = \begin{cases}
	x \quad & \text{ if } x \notin U, \\
	(\phi^{-1} \circ b \circ \phi)(x) \quad & \text{ if } x \in U.
\end{cases}$$
Then by construction, $e(M) \subseteq M$ and $e(p) \notin e(M)$, so we may define 
\begin{equation}\label{eq:sigma}
\sigma: C_n(M) \to C_{n+1}(M)
\end{equation}
by
$$\sigma(m_1,\ldots,m_n) = (e(m_1),\ldots,e(m_n),e(p)).$$
\end{construction}


\subsection{Equivariant configuration spaces}\label{SS:CGS}

In this section, we study equivariant configuration spaces. Our main result (\cref{Cor:CGSDecomp}) will allow us to relate these spaces back to nonequivariant configuration spaces. 

\begin{definition}\label{Def:ConfG-equivariant}
Let $X$ be a $G$-space and let $n\geq 1$. The configuration space of $n$ ordered
points in $X$ is the $G$-space of nonequivariant embeddings
$$ \Conf_n(X) = \Emb(\{ 1,\dots,n \},X). $$
The configuration space of $n$ unordered points in $X$ is the quotient $G$-space
$$ C_n(X) = \Conf_n(X)/\Sigma_n. $$
Let $C(X) = \bigsqcup_{n\geq 1} C_n(X)$ equipped with the disjoint union topology. In $\Conf_n(X)$ and $C_n(X)$, the $G$-action is induced by the action on $X$.
\end{definition}

These are the same definitions as in the nonequivariant case, but now they are
$G$-spaces. In Proposition \ref{Prop:FixptConfDecomp}, we will show that their
fixed points $C_n(X)^G$ split into pieces $C_S^G(X)$ which are described below.

\begin{definition}\label{Def:ConfG}
Let $X$ be a $G$-space and let $S$ be a finite $G$-set of cardinality $n$. The \emph{space of ordered $S$-configurations in $X$}, denoted $\Conf^G_S(X)$, is the space of $G$-equivariant embeddings $S\to X$,
$$\Conf_S^G(X) := \Emb^G(S,X).$$
\end{definition}

This is a $G$-equivariant analogue of the ordinary notion of an ordered configuration space. Indeed, if $G=e$ is the trivial group, we have
$$ \Conf^e_S(X) \cong \Conf_{n}(X^e). $$
Precomposing a $G$-equivariant embedding $S\to X$ by an automorphism of $S$ yields another such embedding. In this way, $\Conf^G_S(X)$ admits an action by the automorphism group $\Aut(S)$. 

\begin{definition} \label{Def:CGS}
The \emph{space of unordered $S$-configurations in $X$} is the quotient
$$ C^G_S(X) := \Conf^G_S(X)/\Aut(S) $$
of $\Conf^G_S(X)$ by the $\Aut(S)$-action.
\end{definition} 

Similarly, this is a $G$-equivariant analogue of the ordinary notion of unordered configuration space. If $G=e$ is the trivial group, then $\Aut(S)\cong\Sigma_{n}$, and we have
$$ C^e_S(X) = \Conf^e_S(X)/\Aut(S) \cong \Conf_{n}(X^e)/\Sigma_{n} = C_{n}(X^e).$$
We can express $C^G_S(X)$ in terms of
nonequivariant configuration spaces of a related space.

\begin{definition}[{\cite[Section I.5]{tom-dieck-transformation-groups}}] \label{Def:X_(H)}
For a $G$-space $X$ and subgroup $H\leq G$, we denote by $X_{(H)}$ the subspace 
$$ X_{(H)} = \{ x\in X : (G_x) = (H) \} \subseteq X $$
consisting of points with stabilizer conjugate to $H$.
\end{definition}

Since the stabilizers of points in a fixed orbit are conjugate, the $G$-action on $X$ restricts to a $G$-action on $X_{(H)}$. The preimage of any point under the projection $X_{(H)}\to X_{(H)}/G$ is a full orbit isomorphic to $G/H$.

\begin{proposition}\label{Prop:CGOrbits}
Given a subgroup $H\subseteq G$ and $k\geq 0$,
the space of unordered $k[G/H]$-configurations in $X$ satisfies 
$$ C^G_{k[G/H]}(X) \cong C_k \left( X_{(H)}/G \right). $$
\end{proposition}
\begin{proof}
A point in $C^G_{k[G/H]}(X)$ is a subset $S\subseteq X$ which is isomorphic as a $G$-set to $k[G/H]$. Each point in $S$ has stabilizer conjugate to $H$, so $S$ is furthermore a subset of $X_{(H)}$. The set $S$ consists of $k$ orbits of type $G/H$, so its image under the projection $X_{(H)}\to X_{(H)}/G$ is a size $k$ subset of $X_{(H)}/G$. This defines a map $$ C^G_{k[G/H]}(X) \to C_k \left( X_{(H)}/G \right), $$
which is a homeomorphism. Indeed, the preimage of any $k$ element subset of $X_{(H)}/G$ is a subset $S\subseteq X_{(H)}\subseteq X$ of the above form, which defines an inverse map.
\end{proof}

The space $X_{(H)}/G$ can also be expressed in terms of $H$-fixed points of $X$ and their $W_GH$-action, as we show in the following lemma.
\begin{lemma}\label{Lem:modW}
For any $G$-space $X$,
$$ X_{(H)}/G \cong \Big( X^H \setminus \bigcup_{K>H} X^K \Big)/W_GH. $$
\end{lemma}
\begin{proof}
There is a $G$-equivariant homeomorphism
\begin{equation}\label{eq:phi} \varphi: X_{(H)} \to G \times_{N_GH} \Big( X^H \setminus
\bigcup_{K>H} X^K \Big) \end{equation}
and hence on orbits we have
\begin{align*}
X_{(H)}/G & 
\cong \Big( X^H \setminus \bigcup_{K>H} X^K \big)/N_GH
\cong \Big( X^H \setminus \bigcup_{K>H} X^K \big)/W_GH
\end{align*}
since $H\leq N_GH$ acts trivially on $X^H$.
\end{proof}

In Corollary \ref{Cor:CGSDecomp} we generalize Proposition \ref{Prop:CGOrbits} to a formula for $C^G_S(X)$ for an arbitrary $G$-set $S$. It follows from the next lemma, whose proof is omitted.

\begin{lemma}\label{Lem:CGProduct}
Let $S$ and $T$ be finite $G$-sets such that no orbit of $S$ is isomorphic to an orbit of $T$. Then,
$$ C^G_{S+T}(X)\cong C^G_S(X)\times C^G_T(X). $$
\end{lemma}
% \begin{proof}
% The assumption guarantees that any embeddings of $S$ and $T$ in $X$ have disjoint images.
% %Indeed, any point in the intersection of any two such embeddings would lie in an orbit contained in both $S$ and $T$.
% It follows that
% \begin{equation}\label{eq:Conf-product} \Conf^G_{S+T}(X) \cong \Conf^G_S(X) \times \Conf^G_T(X).
% \end{equation}
% We similarly have that $\Aut(S+T)\cong\Aut(S)\times\Aut(T)$ (\EvaNoteHl{see}{Isn't this sort of obvious? Because of the
% assumptions on $S,T$, any automorphism has to take $S$ to $S$ and $T$ to $T$,
% and it's determined by that.} e.g.
% \cite[Theorem 12]{CRRM22} for the automorphisms of a $G$-set). Moreover, these
% product decompositions are compatible:
% \begin{align*}
% 	{\Conf^G_S(X) \times \Conf^G_T(X)\over \Aut(S)\x \Aut(T)} &\cong {\Conf^G_S(X)\over \Aut(S)} \times
% {\Conf^G_T(X)\over \Aut(T)}.\qedhere
% \end{align*}
% % Therefore,
% % \begin{align*} 
% % C^G_{S+T}(X) &\cong \Conf^G_{S+T}(X) / \Aut(S+T) \\
% % &\cong \left( \Conf^G_S(X) \times \Conf^G_T(X) \right) / \left( \Aut(S) \times \Aut(T) \right) \\
% % &\cong \left( \Conf^G_S(X) / \Aut(S)  \right) \times \left( \Conf^G_T(X) / \Aut(T) \right) \\ 
% % &= C^G_S(X) \times C^G_T(X)
% % \end{align*} 
% % as desired.
% \end{proof}


Combining \cref{Prop:CGOrbits} and \cref{Lem:CGProduct}, we obtain our desired decomposition formula.

\begin{corollary} \label{Cor:CGSDecomp}
Let $S$ be the finite $G$-set
$$ S = \coprod_{(H)} k_{(H)}[G/H], $$
where $k_{(H)}\geq 0$ for each conjugacy class of subgroups $(H)$. Then,
$$ C^G_S(X) \cong \prod_{(H)} C_{k_{(H)}} (X_{(H)}/G ). $$
\end{corollary}


%The following standard result allows us to choose neighborhoods of orbits of
%points that look like the tangent space at those points.
%\begin{proposition}[``Slice theorem''] % Wikipedia; a historical reference is Palais, ``On the existence of slices for actions of non-compact Lie groups''
%Suppose a finite group $G$ acts smoothly on a smooth manifold $M$. For $x\in M$
%let $G_x$ denote the isotropy subgroup and $T_xM$ the tangent space at $x$. Then
%there exists an invariant neighborhood of $G\cdot x$ equivariantly diffeomorphic
%to $G\x_{G_x} T_xM$.
%\end{proposition}

%\subsection{Expressing fixed configurations in terms of $S$-configurations}

Recall that $C_n(X)$ denotes the $G$-space of nonequivariant embeddings of $n$
points. The next proposition is the main result of this section, which describes
its fixed points in terms of the spaces $C_S^G(X)$ we have been studying above,
in the case that $X$ is a manifold.

\begin{proposition}\label{Prop:FixptConfDecomp}
For a $G$-manifold $M$ and $n\geq 0$, there is a disjoint union decomposition
$$ C_n(M)^G \cong \coprod_{|S|=n} C^G_S(M), $$
where the disjoint union runs over isomorphism classes of finite $G$-sets $S$ with cardinality $n$.
\end{proposition}

\begin{proof}
Let $S = \coprod_{i\in I} G/H_i$ be a $G$-set of cardinality $n$.
Since $C_S^G(M)$ and $C_n(M)^G$ are given compatible quotient topologies, it
suffices to show that $\Conf_S^G(M)$ is open in $\Conf_n(M)^G$.
Let $x\in \Conf_S^G(M)$ be an ordered configuration consisting of orbits
$\mathcal{O}_i= \{ g\cdot x_i \}_{g\in G/H_i}$ for $i\in I$.
By the equivariant slice theorem (see e.g. \cite[Theorem I.2.1.1]{audin}),
there are open neighborhoods $U_{i,g}\ni g\cdot x_i$
and diffeomorphisms $\phi_i:\bigsqcup_g U_{i,g} \cong G\x_{H_i} T_{x_i}M$ with
$\phi_i(g\cdot x_i)= (g, 0)$. Then $U = \prod_{i,g}
U_{i,g}\subseteq M^{\x n}$ is an open neighborhood of $x$.

An element of $\Conf_n^G(M)\ints U$ is an $n$-tuple $(y_{i,g})$ where
$y_{i,g}\in U_{i,g}$, and
we may assume that the $U_{i,g}$'s are disjoint in $M$.
It suffices to show that $y_{i,g}$ has isotropy $H_i$.
Since each $U_{i,g}$ contains a single $y_{i,g}$ and $G$ acts on $\{ y_{i,g} \}$,
an element $g'\in G$ fixes $y_{i,g}$ if and only if it preserves $U_{i,g}\cong
\{ g \}\x T_{x_i}M \subseteq G\x_{H_i} T_{x_i}M$, and this occurs if and only if $g'\in H_i$.
\end{proof}

\begin{remark}
Though Proposition \ref{Prop:FixptConfDecomp} uses the theory of
(finite-dimensional) manifolds, we will investigate an analogous statement for $\rho^\infty$ in future work
and identify the pieces:
\begin{equation}\label{eq:rho-infty} C^G_{G/H}(\rho^\infty)\hteq B(W_GH). \end{equation}
These spaces are interesting because
Guillou and May \cite[Lem. 1.2]{GM17} identify $C_n(\rho^\infty)$ with the $n$-th
space in the $G$-equivariant little $\rho^\infty$-disks operad, which is
well-known to be $G$-equivalent to $B_G\Sigma_n$ (cf. \cite[Pg. 489]{CW91}).
In particular, the splitting here recovers the Lashof--May splitting
\cite[Theorem VII.2.7]{May96} of $(B_G\Sigma_n)^G$.
\end{remark}



%\begin{proof}
%It suffices to show that, given a $G$-set $S = \coprod_i G/H_i$ of cardinality
%$n$, $C_S^G(M)$ is open in $C_n(M)^G$.
%Let $\bar{x}\in C_S^G(M)$ and 
%consider an ordered representative
%$x\in \Conf_S^G(M)\subseteq M^{\x n}$, which we notate $x = (x_{i,j})$ where
%$x_{i,j}$ is in the image of $G/H_i\subseteq S$. We will find an open neighborhood $\bar{U}\subseteq
%C_n^G(M)$ of $\bar{x}$ contained in $C_S^G(M)$.

%By the Slice Theorem (see e.g. \cite[Theorem I.2.1.1]{audin}),
%there exist open neighborhoods $U_{i,j}\ni x_{i,j}$ such that there are
%$G$-equivariant diffeomorphisms
%$\bigcup_j U_{i,j} \cong G\x_{H_i} T_{x_i}M$ for each $i$
%where $x_i=x_{i,j}$ for any $j$ and the image of $U_{i,j}$ is $\{ g_j \} \x
%T_{x_i}M$ with $x_{i,j}$ corresponding to the zero tangent element.
%Since $M$ is Hausdorff, we may shrink these neighborhoods to assume the $U_{i,j}$'s are pairwise disjoint considered as subsets of a single copy of $M$.
%Recall that $C_n^G(M)$ is given the quotient topology via the quotient map
%$q:\Conf_n^G(M)\to C_n^G(M)$. Thus $\bar{U}:= q((\prod U_{i,j})\ints
%\Conf_n^G(M))$ is open.

%Let $\bar{y}\in \bar{U}$; we must show $\bar{y}\in C_S^G(M)$. Let
%$y \in U\subseteq \Conf_n^G(M)$ be a lift to $C_S^G(M)$ and write
%$y = (y_{i,j})$ where $y_{i,j}\in U_{i,j}$. It suffices to show that $y\in
%\Conf_S^G(M)$, and for that it suffices to show that each $y_{i,j}$ has the same
%isotropy type as the corresponding element $x_{i,j}$; that is, $g\cdot
%y_{i,j}=y_{i,j}$ if and only if $g\in H_i$.

%Since $Y = \{ y_{i,j} \}$ is permuted by $G$ and each $U_{i,j}$ contains a
%single element of $Y$, for any $g\in G$ and $y_{i,j}\in Y$ we have that $g$
%fixes $y_{i,j}$ or $g\cdot y_{i,j}\notin U_{i,j}$. Now use the fact that $g$
%preserves $\{ g_j \}\x T_{x_i}M\subseteq G\x_{H_i}T_{x_i}M$ if and only if $g\in
%H_i$.
%\end{proof}

% \begin{lemma}\label{Lemma:ConfGOrbit}
% The space of ordered $G/H$-configurations in $X$ is isomorphic to 
% $$ \Conf^G_{G/H}(X) \cong X^H \setminus \bigcup_{K> H} X^K $$
% where the union runs over all subgroups $K$ containing $H$.
% \end{lemma}
% \begin{proof}
% Recall that the space of $G$-equivariant maps $G/H\to X$ is isomorphic to $X^H$, in which a map $f:G/H\to X$ is determined by a choice of point $f(eH)\in X^H$. We will show that $f$ is an embedding exactly when $f(eH)\notin X^K$ for all subgroups $K$ properly containing $H$. 

% In one direction, suppose that $f(eH)\in X^K$ for some subgroup $K$ such that $H<K$. In this situation, there exists $a\in K$ fixing $f(eH)$ such that $a\notin H$. By the equivariance of $f$, 
% $$ f(eH) = a\cdot f(eH) = f(aH), $$
% but $eH\neq aH$, since $a\notin H$. Therefore, $f$ is not an embedding. 

% Conversely, suppose that such an $f:G/H\to X$ is not an embedding. Without loss of generality, there is a coset $aH\neq eH$ such that $f(aH)=f(eH)$. The image of $f$ is a transitive $G$-set, and is thus isomorphic to $G/K$ for some subgroup $K\leq G$. Thus, $f$ factors as
% $$ \begin{tikzcd}
% G/H \arrow[rr, "f", bend right] \arrow[r, two heads, "p"] & G/K \arrow[r, hook, "f'"] & X
% \end{tikzcd} $$
% for $G$-equivariant maps $p:G/H\to G/K$ and $f':G/K\to X$. Denote the common image under $p$ of $eH$ and $aH$ as 
% $$ \alpha K = p(eH) = p(aH) \in G/K. $$
% It follows from $p(eH)=\alpha K$ that $H\leq\alpha K\alpha\nv$. Furthermore, since $p(aH)=\alpha K$, we have 
% $$ a\nv\alpha K = a\nv p(aH) = p(eH) = \alpha K, $$
% so $a\nv\in \alpha K\alpha\nv$. Since $a\notin H$ by assumption, it follows that $H$ is properly contained in $\alpha K\alpha\nv$. By the equivariance of and $f$ and $f'$, we have for all $k\in K$ that
% $$ (\alpha k\alpha\nv)\cdot f(eH) = (\alpha K\alpha\nv)\cdot f'(\alpha K) = f'(\alpha k\alpha\nv\alpha K) = f'(\alpha K) = f(eH), $$
% so $f(eH)$ is stabilized by the subgroup $\alpha K\alpha\nv$. Therefore, $f(eH)\in X^{\alpha K\alpha\nv}$ where $H<\alpha K\alpha\nv$.

% We conclude the space of $G$-equivariant embeddings $G/H\to X$ is isomorphic to
% $$ X^H \setminus \bigcup_{K>H} X^K, $$
% where the union runs over all subgroups $K$ properly $H$. This space of embeddings is exactly the space $\Conf^G_{[G/H]}(X)$ of ordered $[G/H]$-configurations in $X$.
% \end{proof}

%\ChaseNote{This seems to be fairly well-known; there might be a better reference for this}


\subsection{Equivariant stabilization maps}\label{SS:GConf}
Our goal in this subsection is to define an equivariant version of Construction
\ref{Constr:nonequivariant-stabilization}.

\begin{definition}\label{Def:HStab}
Let $H \leq G$ be a subgroup. A $G$-manifold $M$ is \emph{$H$-stabilizable} if
it is equivariantly homeomorphic to the interior of a $G$-manifold $W$ with boundary containing a point of isotropy group $H$. 
\end{definition}

\begin{construction}\label{Constr:sigmaGH}
Suppose $M$ is an $H$-stabilizable $G$-manifold of dimension $n$, and write $M \isom W^\circ$ where $W$ is a $G$-manifold with boundary containing a point $p \in \partial W$ with isotropy group $H$. By the equivariant slice theorem for manifolds with boundary \cite[Theorem 3.6]{Kan07}, there exists an open subset $U \subseteq W$ containing the orbit $\{ g\cdot p \st g\in G/H \}$ together with a $G$-equivariant diffeomorphism 
%\switchmargin\EvaNote{I don't think the old version $U\to \r_{\geq 0}\oplus V$ (where $U$ had one connected component) was equivariant, since $V$ is a $G$-representation, hence preserved by the $G$-action. I'm not sure I got it right this time, though\dots} %J.D.: I think the proof is fine now (12/19/22)
$$\phi: U\to \r_{\geq 0} \times (G\x_H V) $$ 
sending $g\cdot p$ to $(0, g,0)$, where $\r_{\geq 0}$ has the trivial action and $V \cong T_p(\partial M)$ is the $G$-representation determined by the action of $G$ on the tangent space of $\partial M$ at $p$. Let $b: \r_{\geq 0} \times (G\x_HV) \to \r_{\geq 0} \times (G\x_HV)$ be a smooth, $G$-equivariant bump function which is the identity in the $G\x_H V$ component and sends $(0,g,0)$ to $(1,g,0)$. Let $e: W \to W$ be the $G$-equivariant self-embedding defined by
$$e(x) = \begin{cases}
	x \quad & \text{ if } x \notin U, \\
	(\phi^{-1} \circ b \circ \phi)(x) \quad & \text{ if } x \in U.
\end{cases}$$
Let $g_1,\ldots,g_r$ range over the elements of $G/H$. We define 
\begin{equation}\label{eq:sigmaGH} \sigma_{G/H}: C_n(M) \to C_{n+|G/H|}(M)\end{equation}
by
$$\sigma(m_1,\ldots,m_n) = (e(m_1),\ldots,e(m_n),e(p),g_1e(p),\ldots,g_re(p)).$$
\end{construction}

\begin{remark}
Observe that in \cref{Constr:sigmaGH}, the bump function $\phi$, and hence $\sigma_{G/H}$, preserves isotropy type. Moreover, since $\sigma_{G/H}$ adds a single $G/H$ orbit, it descends to maps
\begin{align*}
C_n(M/G) & \to C_{n+1}(M/G)
\\C_n(M)^G & \to C_{n+|G/H|}(M)^G
\\C_S^G(M) & \to C_{S+[G/H]}^G(M).
\end{align*}
\end{remark}

Putting together Corollary \ref{Cor:CGSDecomp} and Proposition
\ref{Prop:FixptConfDecomp}, we expect equivariant homological stability of $M$ should be
related to nonequivariant homological stability of $M_{(H)}/G$. First (Lemma
\ref{lem:M_(H)/G-open-manifold}) we check that $M_{(H)}/G$ satisfies the
hypotheses of nonequivariant homological stability except for connectedness,
which is an additional assumption and not needed for the definition of
stabilization maps. Then (Lemma
\ref{Lem:sigma-comparison}) we check that the stabilization maps agree. In the
next section (\cref{Thm:CGSHomStab}) we will put these together to prove a form
of equivariant homological stability.

\begin{lemma} \label{lem:M_(H)/G-open-manifold}
If $M$ is an $H$-stabilizable smooth $G$-manifold for $H\leq G$, then
$M_{(H)}/G$ is an open smooth manifold.
\end{lemma}
\begin{proof}
Let $W = W_HG$.
To see that $M_{(H)}/G$ is a manifold, note that $M_{(H)}/G = M_{(H)}/W$ since
$M_{(H)}$ consists entirely of $G/H$-orbit types. Thus $W = \Aut(G/H)$ acts
freely on the manifold $M_{(H)}$, and its acts properly because $W$ is finite. Thus
$M_{(H)}/W$ is a (smooth) manifold.

For noncompactness, first we show noncompactness of $M_{(H)}$. The stabilizability assumption says that there is a
neighborhood of the manifold with boundary $W$ that is diffeomorphic to
$\r_{\geq 0}\x (G\x_H V)$, where $V$ is the tangent space in $\partial W$ of a
point $p\in \partial W$ of isotropy $H$ (see Construction \ref{Constr:sigmaGH}).
The subspace of isotropy $H$ in this neighborhood is $r_{\geq 0}\x (G\x_H V^H)$,
and taking away the boundary, we see that $M_{(H)}$ is missing a limit point,
namely $\{ 0 \}\x (G\x_H \{ 0 \}) = p$. Now passing to orbits $M_{(H)}/G$, we
have a neighborhood $\r_{>0}\x V^H$ with limit point $(0,0)$ that is also not in
$M_{(H)}/G$.
\end{proof}

%\EvaNote{Make it clearer where this is needed}
%\begin{lemma}\label{Lem:UnderlyingMap}
%Suppose $M$ is $H$-stabilizable. Then the underlying map of the equivariant stabilization map \eqref{eq:sigmaGH}
%$$\sigma_{G/H}^e : C_n(M)^e \to C_{n+|G/H|}(M)^e$$
%is homotopic to the $|G/H|$-fold composite $\sigma^{\circ |G/H|}$ of any nonequivariant stabilization map \eqref{eq:sigma}. In particular, it induces the same map as $\sigma^{\circ |G/H|}$ in homology. 
%\end{lemma}

%\begin{proof}
%The homotopy class of the classical stabilization map \eqref{eq:sigma} depends only on the path component of the boundary where the point $p$ is added since the function $e$ is homotopic to the identity (cf. \cite[\S 4]{RW13}). The underlying map $\sigma_{G/H}^e$ factors as the composite of the $r=|G/H|$ nonequivariant stabilizations defined by choosing the points $p, g_1 p, \ldots, g_r p$ in $\partial W$. Since $\partial W$ is assumed to be connected as part of the definition of $H$-stabilizable, we have that $\sigma_{G/H}^e$ is homotopic to the $r$-fold composite of any nonequivariant stabilization map. 
%\end{proof}


\begin{lemma}\label{Lem:sigma-comparison}
Suppose $M$ is $H$-stabilizable.
%and that $M_{(H)}/G$ is an open manifold.
Then the equivariant stabilization map \eqref{eq:sigmaGH} induces a map
$$ \bar{\sigma}_{G/H}: C_n(M_{(H)}/G)\to C_{n+1}(M_{(H)}/G) $$
that agrees with the nonequivariant stabilization map \eqref{eq:sigma}. 
\end{lemma}
Note that Lemma \ref{lem:M_(H)/G-open-manifold} is used to guarantee the map
\eqref{eq:sigma} exists.
\begin{proof}
Since $\sigma_{G/H}:C_n(M)\to C_{n+|G/H|}(M)$ preserves isotropy type and adds a $G/H$ orbit, it induces a map $\bar{\sigma}_{G/H}^e:C_n(M_{(H)}/G)\to C_{n+1}(M_{(H)}/G)$, and the equivariant bump function $b$ induces a function $\bar{b}$ on $M_{(H)}/G$. First note from Lemma \ref{lem:M_(H)/G-open-manifold} that $M_{(H)}/G$ is a manifold.
To check $\bar{\sigma}_{G/H}$ agrees with the classical stabilization map, we must check that $\bar{b}$ is a bump function. 
It is clear that it has compact support. To check it is smooth, first note that the restriction of the smooth function $b$ to $M_{(H)}$ is smooth. To check it is smooth after taking orbits, recall that every orbit in $M_{(H)}$ has a neighborhood diffeomorphic to $G\x_H \r^n$, and taking orbits simply collapses the $|G/H|$ Euclidean components to one Euclidean component.
\end{proof}




\section{Equivariant homological stability}\label{Sec:Stab}

In this section, we prove our main theorems, \cref{MT:BredonFails} and \cref{MT:Stability}. Our proof is by reduction to the classical stability result of McDuff--Segal (\cref{Thm:HomStab}), which we recall in \cref{SS:MS}. We also discuss a weaker notion of homological stability (\cref{Cor:WeakStability}) in terms of finite generation; the relationship between finite generation stability and sequential stability is made precise in \cref{Prop:AlgStability}. Using nonequivariant homological stability, we prove \cref{MT:CSGStability} (\cref{Thm:CGSHomStab}) that the equivariant stabilization maps eventually induce isomorphisms in the homology of $S$-configuration spaces in \cref{SS:CGSStability}. We then use our decomposition of the $G$-fixed points of configuration spaces into $S$-configuration spaces, along with our Bredon-to-geometric-fixed-point reduction result (\cref{Lemma:Bredon-to-PhiG}), to prove Bredon homological stability in \cref{SS:BredonStab}.

\subsection{Nonequivariant homological stability}\label{SS:MS}

\begin{theorem}[{McDuff--Segal, as stated in \cite[Thm. 1.2]{Pal18} with $X=*$}]\label{Thm:HomStab}
Assume $M$ is an open, connected manifold with $\dim(M) \geq 2$. The stabilization map $\sigma:C_n(M) \to C_{n+1}(M)$ induces an isomorphism on integral homology in degrees $* \leq \frac{n}{2}.$
\end{theorem}

Our goal for the rest of this subsection is to relate homological stability with respect to iterated applications of $\sigma$ with finite generation over $\z[\sigma_*]$. Since this does not depend on the fact that the groups that stabilize are homology groups, we take an abstract approach, presenting the translation for arbitrary sequences of modules that stabilize.

Fix a commutative noetherian ring $R$ and let 
$$ A_0 \xrightarrow\sigma A_1 \xrightarrow\sigma A_2 \xrightarrow\sigma \cdots $$
be a sequence of $R$-modules $A_n$ for $n\geq 0$ connected by maps $\sigma_n:A_n\to A_{n+1}$. Write 
$$ A = \bigoplus_{n\geq 0} A_n. $$
The collection of maps $\sigma_n$ assemble into an endomorphism $\sigma:A\to A$, and this map endows $A$ with the structure of a $R[\sigma]$-module. The next proposition relates finiteness conditions on $A$ and $\sigma$ to a stability condition on the sequence.

\begin{proposition} \label{Prop:AlgStability}
With $R$, $A$ and $\sigma$ as above, the following are equivalent:
\begin{enumerate}
\item $A_n$ is a finitely generated $R$-module for each $n$, and the sequence stabilizes: that is, there exists $N\geq 0$ such that for all $n\geq N$, $\sigma:A_n\to A_{n+1}$ is an isomorphism; and
\item $A$ is finitely generated as an $R[\sigma]$-module.
\end{enumerate}
\end{proposition}

\begin{proof}
To see (1) $\Rightarrow$ (2), let $f_n:R^{r_n}\to A_n$ be a surjection. Define a map
$$ f:R[\sigma]^{r_0} \oplus \cdots \oplus R[\sigma]^{r_N} \to A $$
which maps the generators of $R[\sigma]^{r_n}$ to $A_n\subseteq A$ according to $f_n$ and extends $R[\sigma]$-linearly. To see that $f$ is surjective, note that by construction it surjects onto $\bigoplus_{n=0}^N A_n$, and it surjects onto $\bigoplus_{n=N}^\infty A_n \cong A_N \otimes \Z[\sigma]$ because it is a $\sigma$-linear map that surjects onto $A_N$.
%Let $a\in A_n$ for some $n\geq 0$. If $n\leq N$, then $a$ is in the image of $f$ by construction, so let $n=N+k$ for some $k>0$. By the stability assumption,
%$$ a=\sigma_{N+k-1}\circ\cdots\circ\sigma_N(a') $$
%for some $a'\in A_N$. The element $a'$ is in the image of $f$, so it has a preimage $r$. By $R[\sigma]$-linearity, we have
%$$ f(\sigma^kb) = \sigma^k\cdot f(b) = \sigma_{N+k-1}\circ\cdots\circ\sigma_N(f(b)) = a, $$
%so $a$ is in the image of $f$. Since $A$ admits a surjection from a free $R[\sigma]$-module of finite rank, $A$ is a finitely generated $R[\sigma]$-module.

To see (2) $\Rightarrow$ (1), let $f:R[\sigma]^r\to A$ be a surjection from a free $R[\sigma]$-module of finite rank $r$. Since $f(\sigma R[\sigma]^r)\subseteq \sigma A$, $f$ descends to a surjective map on quotients 
$$ \bar{f}: R^r \cong R[\sigma]^r/\sigma R[\sigma]^r \to A/\sigma A = \coker\sigma$$
which exhibits $\coker\sigma$ as a finitely generated $R$-module. 
% (To see that $\bar{f}$ is surjective, note that it fits into the commutative diagram
% \[
% \begin{tikzcd}
% R[\sigma]^r \arrow[r,twoheadrightarrow,"f"] \arrow[d,twoheadrightarrow] & A \arrow[d,twoheadrightarrow] \\
% R[\sigma^r]/\sigma \arrow[r,dashed,"\bar{f}"] & A/\sigma
% \end{tikzcd}
% \]
% where each solid arrow is surjective.)
Note that 
$$ \coker\sigma \cong \bigoplus_{n\geq 0} (\coker\sigma_n). $$
That $\coker\sigma$ is finitely generated thus implies that $(\coker\sigma)_n=0$ for sufficiently large $n$. It follows that the sequence is eventually surjective, that is, there exists a $M$ such that $\sigma_n$ is surjective for $n\geq M$. If we let
$$ K_i = \ker(A_M \to A_i), $$
for $i\geq M$, we obtain an increasing chain $K_M\subseteq K_{M+1}\subseteq\cdots$ of submodules of $A_M$. As a finitely generated module over the noetherian ring $R$, $A_M$ is noetherian. Therefore, there must exist an $N\geq 0$ at which point this increasing chain stabilizes, that is $K_N=K_{N+1}=\cdots$. It follows that for $n\geq N$, $\sigma_n$ is an isomorphism.

Now we show that $A_n$ is finitely generated over $R$ by induction on $n$. Since
$\coker(\sigma)|_{A_0} = A_0$, we must have that $A_0$ is finitely generated.
Now assume $n\geq 1$ and assume inductively that $A_{n-1}$ is finitely
generated. Consider the exact sequence
$$ A_{n-1}\too{\sigma} A_n \too{\sigma} \coker(\sigma)|_{A_n}\to 0. $$
Since $A_{n-1}$ and $\coker(\sigma)|_{A_n}$ are finitely generated by
assumption, we have that $A_n$ is also finitely generated.
\end{proof}

\begin{corollary}\label{Cor:WeakStability}
Let $M$ be a stabilizable open, connected manifold with $\dim(M) \geq 2$. If $H_d(C_k(M))$ is finitely generated as an abelian group for all $d \geq 0$, then
$H_*(\bigsqcup_{k \geq 0} C_k(M))$ is finitely generated as a $\z[\sigma_*]$-module. 
\end{corollary}

We note that the finite generation of $H_d(C_k(M))$ in \cref{Cor:WeakStability} follows from the finite generation of $H_d(M)$. First we need a lemma.

\begin{lemma} \label{lem:homology-orbits-finite}
Suppose $X$ has a free $G$-action, and $H_d(X;\Z)$ is a finitely generated
abelian group for all $d$. Then $H_d(X/G;\Z)$ is finitely generated.
\end{lemma}
\begin{proof}
There is a fibration $X\to EG\x_G X \to BG$, and
since the action of $G$ on $X$ is free, we have $EG\x_G X\hteq
X/G$. There is a Serre spectral sequence with local coefficients
$$ E_2^{*,*}= H_*(BG; H_*(X;\Z)) \Rightarrow H_*(X/G). $$
The $E_2$-page is the cohomology of the complex $C_*(EG;\Z)\otimes_{\Z[G]}
H_*(X;\Z)$ which is finitely generated in each degree by assumption.
Thus the $E_2$-term is finitely generated in each degree, and so is the
target.
\end{proof}

\begin{proposition}\label{Prop:MtoCnM}
If $H_d(M)$ is finitely generated for all $d \geq 0$, then so is $H_d(C_n(M))$ for all $d,n \geq 0$. 
\end{proposition}
\begin{proof}
By Lemma \ref{lem:homology-orbits-finite}, it suffices to show that
$H_d(\Conf_n(M);\Z)$ is finitely generated. Since $\Conf_n(M)$ is a manifold,
its homology is Poincar\'e dual to cohomology with compact support
$H^d_c(\Conf_n(M);\mathcal{O})$ with coefficients in the orientation sheaf $\mathcal{O}$ (see
\cite[Example 2.9]{bredon-sheaf-theory}).
To study this, we apply \cite[Theorem 5.18]{petersen-generalized}, which gives a
spectral sequence that in our case specializes to
$$ E_1^{p,q} = \bigoplus_\attop{T\in J_{U_0}\\|T| = n-p}\bigoplus_{i+j=p+q} \til{H}^i(J(T);
H^j_c(\overline{M(T)}; \mathcal{O})) \implies H_c^{p+q}(\Conf_n(M); \mathcal{O}) $$
where $J_{U_0}$ is a subset of the set of partitions of $\{ 1,\dots,n \}$,
$J(T)$ is a finite CW complex, and $\overline{M(T)}$ is a closed submanifold of 
$$ \{ (x_1,..,x_n)\in M^n \st x_i=x_j \iff \text{$x_i = x_j$ if $i$ and $j$ are in the
same block of $T$} \}, $$
which is itself a closed submanifold of $M^n$. Closed manifolds have cohomology
that is finitely generated in each degree. Since the $E_1$ page of the spectral
sequence is finitely generated for each $(p,q)$ and $p,q\geq 0$ we have that the
target is finitely generated in each degree.
\end{proof}

\subsection{Instability for adding a trivial orbit}

If $M$ is $G$-stabilizable, then we may add a fixed point to obtain a $G$-equivariant stabilization map
$$\sigma_{G/G}: C_n(M) \to C_{n+1}(M)$$
whose underlying map
$$\sigma_{G/G}^e = \sigma : C_n(M) \to C_{n+1}(M)$$
is the classical stabilization map. In this section, we show that the Bredon homology of $C(M)$ cannot be stable with respect to this stabilization map alone. For simplicity of exposition, we consider the degree 0 case where $G=C_p$.

% We may therefore apply \cref{Lemma:Bredon-to-PhiG} to identify
% $$\coker(H_d^G(C_n(M);\mfz) \xrightarrow{(\sigma_{G/G}^e)_*} H_d^G(C_{n+1}(M);\mfz))$$
% with 
% $$\coker(\pi_d(\Phi^G(H\mfz \wedge C_n(M))) \xrightarrow{(\sigma_{G/G}^e)_*} \pi_d(\Phi^G(H\mfz \wedge C_{n+1}(M))))$$
% for $n \gg 0$. Using the properties of geometric fixed points established in \cref{Sec:Eqvt}, we have
% $$\pi_d(\Phi^G(H\mfz \wedge C_k(M))) \cong \bigoplus_{i \in I} H_{d-n_i}( C_k(M)^G; \f_p)$$
% assuming $G$ is a $p$-group (if $G$ is not a $p$-group, then this group is zero), where $I$ is a finite set of integers with $0 \leq n_i \leq d$ for each $i \in I$. This reduces the question of Bredon homological stability to the question of homological stability for the fixed points $C_k(M)^G$. 

% \begin{proposition} 
% Suppose $G$ is a finite group with a subgroup of index $k>1$, and write $\sigma = \sigma_{G/G}^e$. Then
% $$(\sigma^{\circ k})_* : H_0(C_n(M)^G;\z) \to H_0(C_{n+k}(M)^G;\z)$$
% is not surjective for any $n \geq 0$.
% \end{proposition}

% \begin{proof}
% We have
% $$C_n(M)^G \cong \coprod_{|S|=n} C_S^G(M), \quad C_{n+k}(M)^G \cong \coprod_{|T|=n+k} C_T^G(M).$$
% The map $\sigma^{\circ k}$ sends an $S$-configuration to an $(S + k[G/G])$-configuration. If $|G/H| = k$ then 
% $$ T = \lfloor{\textstyle{n+k\over k}}\rfloor [G/H] + (n+k-k\lfloor{\textstyle{n+k\over k}}\rfloor) [G/G] $$
% is a $G$-set of cardinality $n+k$ that contains less than $k$ copies of the trivial orbit $[G/G]$, and therefore the corresponding summand is not in the image of $\sigma^{\circ k}$.
% The result follows. 
% \end{proof}

\begin{proposition} \label{Prop:H0Unstable}
For $G=C_p$ and $M$ an open $G$- and $e$-stabilizable manifold, the $\z[(\sigma_{G/G})_*]$-module $H_0^G(C(M);\mfz)$ is not finitely generated. 
\end{proposition}
\begin{proof}
The cellular chain complex $C^G_*(C(M))$ computing Bredon homology (cf. \cref{SS:Bredon}) is given by
\begin{align*}
C^G_*(C(M)) &= \Big( \sum_{G/H} C_*(C(M)^H)\otimes\mfz(G/H) \Big)/\sim\\
&= \Big( C_*(C(M)^e) \otimes \mfz(C_p/e) \oplus C_*(C(M)^G) \otimes \mfz(C_p/C_p) \Big) /\sim,
\end{align*}
% where the relation $\sim$ is $f^*x\otimes y\sim x\otimes f_*y$ for $f:G/H\to G/K$ in the orbit category, $x\in C_*(C(M)^K)$, and $y\in \mfz(G/H)$. 
In the case $G=C_p$, the orbit category is generated by the projection $r:C_p/e\to C_p/C_p$ and automorphism $\gamma:C_p/e\to C_p/e$ which corresponds to the action of the generator of $C_p$. Therefore, in $C^G_*(C(M))$, we have the relations:
\begin{itemize}
\item $\gamma\cdot x\sim x$ for $x\in C_*(C(M)^e)$ and $\gamma\in C_p$ the generator of $C_p$, and 
\item $r^*x\sim px$ for $x\in C_*(C(M)^G)$, where $r^*:C(M)^G\to C(M)^e$ is the inclusion of fixed points.
\end{itemize} 
From \cref{Prop:FixptConfDecomp}, we have that the components of $C(M)^G$ are indexed by finite $G$-sets, so we have
$$ H_0(C(M)^{C_p}) \cong \ZZ[\sigma]\{ x_0, x_1, x_2, \ldots \}, $$
where $x_i$ represents the component corresponding to the $G$-set $i\cdot C_p/e$. The components of $C(M)$ are indexed by natural numbers, so we have 
$$ H_0(C(M)^e) \cong \ZZ[\sigma]\{ y_0 \}. $$
We conclude that $H^G_0(C(M))$ is given by
$$ H^G_0(C(M)) \cong \frac{ \ZZ[\sigma]\{ x_0,x_1,x_2,\cdots, y_0 \} }{ px_i = \sigma^i y_0}, $$
which is not finitely generated as a $\ZZ[\sigma]$-module.
\end{proof}

\subsection{Homological stability for $S$-configurations}\label{SS:CGSStability}

%\subsection{Homological stability for $C_S^G(M)$}

As above, let $M$ be a $G$-manifold. In this section, we fix a conjugacy class of subgroups $(H)$ and study stabilization maps for $[G/H]$ orbits. We will combine Corollary \ref{Cor:CGSDecomp} with classical homological stability applied to the nonequivariant space $M_{(H)}/G$, to obtain homological stability for $C_S^G(M)$. 


\begin{theorem}\label{Thm:CGSHomStab}
Let $M$ be a $G$-manifold and $S$ a finite $G$-set containing $k$
copies of the orbit
$[G/H]$ for a conjugacy class of subgroups $(H)$. If $M$ is $H$-stabilizable and
$M_{(H)}/G$ is connected,
%and $M_{(H)}/G$ is an open connected manifold,
then the $(H)$-stabilization map of \eqref{eq:sigmaGH} induces an isomorphism
$$ (\sigma_{G/H})_*: H_d(C^G_S(M)) \to H_d(C^G_{S+[G/H]}(M)) $$
in integral homology in degrees $d\leq k/2$.
\end{theorem}

\begin{proof}
Decompose the finite $G$-set $S$ as 
$$ S \cong k[G/H] \sqcup S', $$
where $S'$ is a finite $G$-set with no orbits isomorphic to $[G/H]$. \cref{Cor:CGSDecomp} allows us to decompose the space of unordered $S$-configurations in $M$ as 
$$ C^G_S(M) \cong C^G_{k[G/H]}(M) \times C^G_{S'} \cong C_k(M_{(H)}/G) \times C^G_{S'}(M). $$
By \cref{Cor:CGSDecomp}, we have 
\begin{align*}
H_*(C^G_S(M)) &\cong H_*(C_k(M_{(H)}/G)\times C^G_{S'}(M)), \\
H_*(C^G_{S+[G/H]}(M)) &\cong H_*(C_{k+1}(M_{(H)}/G)\times C^G_{S'}(M)).
\end{align*}

By naturality, we obtain a map of K\"unneth short exact sequences
%\ChaseNote{It's better but I still don't love it.}
% $$ \tiny \begin{tikzcd}
% 0 \arrow[r] & \displaystyle\bigoplus_{i+j=d-1} \Tor_1\left(H_i(C_k(M_{(H)}/G)), H_j(C^G_{S'}(M))\right) \arrow[r] \arrow[d] & H_d(C^G_S(M)) \arrow[r] \arrow[d, "(\sigma_{(H)})_*"] & \displaystyle\bigoplus_{i+j=d} H_i(C_k(M_{(H)}/G))\otimes H_j(C^G_{S'}(M)) \arrow[r] \arrow[d] & 0 \\
% 0 \arrow[r] & \displaystyle\bigoplus_{i+j=d-1} \Tor_1\left(H_i(C_{k+1}(M_{(H)}/G)), H_j(C^G_{S'}(M))\right) \arrow[r] & H_d(C^G_{S+[G/H]}(M)) \arrow[r] & \displaystyle\bigoplus_{i+j=d} H_i(C_{k+1}(M_{(H)}/G))\otimes H_j(C^G_{S'}(M)) \arrow[r] & 0
% \end{tikzcd} $$
$$ \adjustbox{scale=0.85,center}{
\begin{tikzcd}
\displaystyle\bigoplus_{i+j=d-1} \Tor_1\left(H_i(C_k(M_{(H)}/G)), H_j(C^G_{S'}(M))\right) \arrow[r] \arrow[d]
& \displaystyle\bigoplus_{i+j=d-1} \Tor_1\left(H_i(C_{k+1}(M_{(H)}/G)), H_j(C^G_{S'}(M))\right) \arrow[d] \\
H_d(C^G_S(M)) \arrow[r, "(\sigma_{G/H})_*"] \arrow[d]
& H_d(C^G_{S+[G/H]}(M)) \arrow[d] \\
\displaystyle\bigoplus_{i+j=d} H_i(C_k(M_{(H)}/G))\otimes H_j(C^G_{S'}(M)) \arrow[r] 
& \displaystyle\bigoplus_{i+j=d} H_i(C_{k+1}(M_{(H)}/G))\otimes H_j(C^G_{S'}(M))
\end{tikzcd}
}
$$
% $$ \begin{tikzcd}
% E^2_{p,q} = \displaystyle\bigoplus_{i+j=q} \Tor_p\left(H_i(C_k(M_{(H)}/G)), H_j(C^G_{S'}(M))\right) \arrow[Rightarrow, r] \arrow[d] & H_{p+q}(C^G_S(M)) \arrow[d, "(\sigma_{(H)})_*"] \\
% E^2_{p,q} = \displaystyle\bigoplus_{i+j=q} \Tor_p\left(H_i(C_{k+1}(M_{(H)}/G)), H_j(C^G_{S'}(M))\right) \arrow[Rightarrow, r] & H_{p+q}(C^G_{S+[G/H]}(M)).
% \end{tikzcd} $$
By \cref{Lem:sigma-comparison}, the top and bottom horizontal arrows in the above diagram are given by the classical stabilization maps $\sigma:C_k(M_{(H)}/G)\to C_{k+1}(M_{(H)}/G)$
for the manifold $M_{(H)}/G$, and so \cref{Thm:HomStab} implies that these maps are isomorphisms for $d\leq k/2$.
(Note that $M_{(H)}/G$ is not empty because we have assumed that there are $k$
points of isotropy type $(H)$.)
It follows from the five lemma that $(\sigma_{G/H})_*$ is an isomorphism in the range $d\leq k/2$.
\end{proof}

\subsection{Bredon homological stability}\label{SS:BredonStab}

We now return to the situation of ordinary configuration spaces.

%Let $M$ be an
%equivariant manifold and let
%$$ C(M) = \bigsqcup_{n\geq 0} C_n(M). $$

For $K\leq G$, let 
\[ P_K = \Z[\sigma_{K/H} : H\subseteq K] . \]
Clearly, $P_K$ acts on $K$-Bredon cohomology $H^K_*(C(M))$, but we can also
consider the action of $P_G$ on $H^K_*(C(M))$ via the restriction
map $\res^G_K: P_G\to P_K$, which is defined as follows.
Given $H,K\leq G$, the double coset formula gives an isomorphism of
$K$-sets
$$ G/H\ \ \cong \bigsqcup_{KgH\in K\backslash G/H} K/(K\ints gHg^{-1}),$$
and so
\begin{equation}\label{eq:res}\res^G_K(\sigma_{G/H})\ \  = \prod_{KgH\in K\backslash
G/H}\sigma_{K/(K\ints gHg^{-1})}.\end{equation}

\begin{lemma} \label{lem:RK-to-RG}
Suppose $G$ is a Dedekind group (i.e., every subgroup is normal).
Then for all $K\leq G$, $P_K$ is finitely generated as a $P_G$-module, where the action is via the restriction map $\res^G_K:P_G\to P_K$.
\end{lemma}
\begin{proof}
As $P_K$ is a finitely generated algebra, it suffices to show that $P_K$ is
integral over $P_G$. We will show that there is a power of every generator
$\sigma_{K/H}$ (for $K\leq H$) in the image of $\res^G_K$, which implies the same for every
monomial. Combining the Dedekind assumption with \eqref{eq:res} shows that
$\res^G_K(\sigma_{G/H}) = \sigma_{K/H}^{[G:K]}$. 
\end{proof}

\begin{remark} \label{rmk:abelian}
The ``Dedekind'' hypothesis is needed to ensure that $\res^G_K(\sigma_{G/H}) = \sigma_{K/H}^{[G:K]}$ in the last line of the previous proof. It is known \cite{Ded97} that a finite group is Dedekind if and only if it is abelian or 
isomorphic to $Q_8\x A$ where $A$ is a finite abelian group whose 2-component 
has the form $(\Z/2)^r$
for $r\geq 0$.

We remark that the conclusion of \cref{lem:RK-to-RG} holds for a more general class of nonabelian groups than those mentioned in the previous paragraph, but an explicit characterization is not known. For example, it can be checked that it holds for dihedral groups of order $2n$ for $n$ squarefree.
\end{remark}

\begin{lemma} \label{lem:PhiG-generated}
Assume $M$ is $H$-stabilizable and $M_{(H)}/G$ is connected for every $H\leq G$. For every $d$ there exists $N_d$ such that
$$ \pi_d(\Phi^G(H\underline{\Z}\sm C(M))) \isom P_G\cdot
\pi_d(\Phi^G(H\underline{\Z}\sm C_{\leq N_d}(M))). $$
Moreover, if $H_i(M_{(H)}/G)$ is a finitely generated abelian group for all $i\leq d$ and $(H)\leq G$, then
$\pi_d(\Phi^G(H\underline{\Z}\sm C_{\leq N_d}(M)))$ is a finitely generated abelian group.
\end{lemma}
\begin{proof}
Using the splitting in \cref{Cor:CGSDecomp} we have
\[
\pi_d(\Phi^G(H\underline{\Z}\sm C(M))) \isom \pi_d(\Phi^GH\underline{\Z}\sm C(M)^G)
\isom \bigoplus_{G\text{-sets }S} \pi_d(\Phi^G(H\underline{\Z})\sm C_S^G(M)).
\]
By Proposition \ref{Prop:Green}, there is a finite collection of nonnegative integers $n_i$ and finitely generated abelian groups $A_i$ such that 
\[ \pi_d(\Phi^G(H\underline{\Z})\sm C_S^G(M))\isom \bigoplus_{i} H_{d-n_i}(C_S^G(M); A_i). \]

For fixed $i$ and $H$, Theorem \ref{Thm:CGSHomStab} says that there exists
$N_{d,i,H}$ such that
$$H_{d-n_i}(C_S^G(M); A_i) = \sigma_{G/H}\cdot H_{d-n_i}(C_{S-[G/H]}^G(M);
A_i) $$
whenever $S$ contains at least $N_{d,i,H}$ copies of $G/H$. For fixed $i$, there
are finitely many $G$-sets $S$ such that $H_{d-n_i}(C_S^G(M);A_i)$ is not in the
image of $\sigma_{G/H}$ for some $H$. Let $N_{d,i}$ be the maximum cardinality
of the exceptional $G$-sets, and
let $N_d = \max_i\{ N_{d,i} \}$. Then for every $S$ and $n_i$, we have
\begin{align*}
H_{d-n_i}(C_S^G(M); A_i) &\subseteq P_G\cdot \bigoplus_{|S'| \leq N_d} H_{d-n_i}(C_{S'}^G(M); A_i).
\end{align*}
Summing over $i$, we obtain $\pi_d(\Phi^G(H\underline{\Z}\sm C^G_S(M)))\subseteq
P_G \cdot \pi_d(\Phi^G(H\underline{\Z}\sm C_{\leq N_d}(M)))$.

This shows the first statement. For the second statement, combine Corollary \ref{Cor:CGSDecomp} and Proposition \ref{Prop:MtoCnM}, which holds for homology with $A_i$-coefficients because $A_i$ is a finitely generated abelian group.
\end{proof}


\begin{proposition}\label{Prop:BredonHomStab}
Let $G$ be a Dedekind group. For $K\leq G$, let $i_K^*$ denote restriction from
$G$-spaces to $K$-spaces.
Assume that for all $H\leq K \leq G$, $M$ is
$H$-stabilizable, $(i_K^*M)_{(H)}/K$ is connected,
and $H_i(i_K^*(M)_{(H)}/K)$ is a finitely generated abelian group for all $i\leq d+1$.
Then $H_d^K(C(M);\underline{\Z})$ is finitely generated over $P_G$ for all $K\leq G$.
\end{proposition}
\begin{proof}
We will apply Lemma \ref{Lemma:Bredon-to-PhiG} with the Serre class of finitely generated $P_G$-modules. It suffices to check the hypothesis, namely that $\pi_i(\Phi^K(H\underline{\Z}\sm C(M)))$ is finitely generated over $P_G$ for all $K\leq G$ and $i\leq d+1$. 
We will apply Lemma \ref{lem:PhiG-generated} to $i^*_K(M)$. Checking the
hypotheses, first note that if $M$ is $H$-stabilizable as a $G$-manifold, then
$i^*_K(M)$ is $H$-stabilizable.
Thus $\pi_i(\Phi^K(H\underline{\Z}\sm C(i^*_K(M))))=\pi_i(\Phi^K(H\underline{\Z}\sm C(M)))$ is finitely generated over $P_K$. Now use Lemma \ref{lem:RK-to-RG} to get finite generation over $P_G$.
\end{proof}

%\begin{lemma} \label{Lem:finitely generated}
%Assume $M$ is $H$-stabilizable and $M_{(H)}$ is noncompact for
%all $(H)\subseteq G$.
%Let $G$ be a finite abelian group and let $S$ be a finite $G$-set. Then there is
%a decomposition
%\begin{equation}\label{eq:Hcal-decomp}\pi_d(\Phi^G(H\underline{\Z}) \sm C_S^G(M)) = \bigoplus_{n_i}
%H_{d-n_i}(C_S^G(M); \Z/p)\end{equation}
%for a collection $\{ n_i \}$ of nonnegative indices that depends only on $G$,
%and this is a finitely generated abelian group.
%\end{lemma}

%\begin{proof}
%If $G$ is not a $p$-group then $\Phi^G(H\underline{\Z})=0$ by
%\cite[Proposition 11]{Kri20}.
%\EvaNote{Proposition \ref{Prop:Green} and \ref{Prop:GFPVanish} are the more
%general versions of the things I'm quoting from S. Kriz,
%but I don't know how to show more generally that there are finitely many $n_i$'s
%that are $\leq d$ for a fixed $d$.}
%By \cite[Proposition 9]{Kri20} it further suffices to reduce to the case of an
%elementary abelian $p$-group $G = (\Z/p)^r$.

%By Proposition \ref{Prop:Green}, there is a collection $\{ n_i \}_{i\in I}$ of
%nonnegative indices $n_i$ such that $H\underline{\Z} = \bigvee_{i\in I}
%\Sigma^{n_i} H\underline{\Z/p}$. In the case $G = (\Z/p)^r$, the collection
%$\{ n_i \}$ can be read off Kriz's computation \cite[Theorem 3]{Kri20} 
%of $\pi_*\Phi^GH\underline{\Z}$; in particular, for every $d\geq 0$, there are
%finitely many summands $\Sigma^{n_i}H\underline{\Z/p}$ with $n_i \leq d$. Since
%$C_S^G(M)$ is a space, its homology is concentrated in nonnegative degrees, so we
%have
%$$ \pi_d(\Phi^G(H\underline{\Z}) \sm C_S^G(M)) = \bigoplus_{i\in I}
%H_{d-n_i}(C_S^G(M); \Z/p) = \bigoplus_\attop{i\in I\\n_i\leq d}
%H_{d-n_i}(C_S^G(M); \Z/p) $$
%which shows there are finitely many summands. 
%By the product decomposition in Proposition \ref{Cor:CGSDecomp} and the
%K\"unneth isomorphism, it suffices to show that $H_k(C_n(M_{(H)}/G))$
%is finitely generated for all subgroups $H$ and all $k,n$.

%\eva{To use the fact from \cite[p.92]{McD75} we discussed, $M_{(H)}/G$ would
%need to be tame, and I'm having trouble showing that using the assumption that
%$M$ is $H$-tame. The closure of a bounded open manifold need not be a manifold;
%if $W$ is the compact manifold containing $M$, I'd like to use
%$W_{(H)}/G\contains M_{(H)}/G$ but the former might not be closed (e.g. $H =
%C_2$ and $M = $ the regular representation of $C_2$).
%}
%Since $W'$ is covered by finitely many closed contractible sets, an inductive
%argument using the Mayer-Vietoris sequence shows that $H_k(C_n(W'/G);\Z/p)$ (and
%hence $H_k(C_n(M_{(H)}/G);\Z/p)$) is finite for all $k,n$.
%By \cite[p.91]{McD75}, under our assumptions on $M$ we have an isomorphism
%$H_d(C_n(M)) \isom H_d(\Gamma_n(M))$ for large $n$, where $\Gamma_n(M)$ is the space of
%degree-$n$ cross-sections of a certain fiber bundle associated to $M$, and the
%proof of \cite[Theorem 4.5]{McD75} shows 
%\end{proof}

%\begin{theorem}\label{Thm:BredonHomStab}
%Let $d \geq 0$ be an integer and let $\uR$ be a Green functor with $\uR(G/H)$ a finitely generated abelian group for all $H \leq G$. 
%Let $M$ be an equivariantly stabilizable $G$-manifold with $M_{(H)/G}$ noncompact, connected and $H_i(M_{(H)}/G)$ finitely generated as an abelian group for all $0 \leq i \leq d$ and $(H) \leq G$. 
%Then the $\z[(\sigma_{G/H})_*: (H) \leq G]$-module $H_d^G(C(M);\mfz)$ is finitely generated.
%\end{theorem}

%\begin{proof}
%For ease of notation write $\H_d(X)=\pi_d(\Phi^G(H\underline{R})\sm X)$ for a
%space $X$, so
%$$ \pi_d(\Phi^G(H\underline{R}\sm C(M))) = \pi_d(\Phi^G(H\underline{R})\sm
%C(M)^G) = \H_d(C(M)^G). $$
%By the splitting in Proposition \ref{Prop:FixptConfDecomp}, we have a
%decomposition
%\begin{equation}\label{eq:H-decomp} \H_d(C(M)^G) =\bigoplus_{n\geq 0}
%\H_d(C_n(M)^G) = \bigoplus_{G\text{-sets }S} \H_d(C^G_S(M)).
%\end{equation}

%Order the conjugacy classes of subgroups of $G$ so the collection
%$\{ \sigma_{G/H} \}_{(H)\subseteq G}$ is written $\{\sigma_1,\dots,\sigma_k\}$.
%Identify the $G$-set $S$ with the tuple $(S_1,\dots,S_k)$ where $S_i$ denotes
%the number of copies of $[G/H_i]$ in $S$. 
%So \eqref{eq:H-decomp} becomes
%$$ \H_d(C(M)^G) = \bigoplus_{S = (S_1,\dots,S_k)\in \N^k} \H_d(C^G_S(M)). $$
%In this notation, $\sigma_i$ restricts to a map $\H_d(C^G_S(M))\to
%\H_d(C^G_{S+e_i}(M))$ where $e_i$ is the $k$-tuple with $1$ in the $i^{th}$
%position and 0 elsewhere.

%For $n\geq 0$ let
%$$ A_n = \bigoplus_{S = (S_1,\dots,S_{k-1},n)} \H_d(C_S^G(M)). $$
%Then $\bigoplus_{n\geq 0} A_n = \H_d(C(M)^G)$.
%We will show by induction on $k$ that $A_n$ is finitely generated over
%$\Z[\sigma_1,\dots,\sigma_{k-1}]$. The base case $k=1$ is \cref{Lem:finitely generated}.
%%\EvaNote{As of now, the hypotheses of Lemma \ref{Lem:finitely generated} need to
%%be added to this theorem; hoping to do better than this.}
%For the inductive step, assume that $A_{n,m}:= \bigoplus_{S =
%(S_1,\dots,S_{k-2},m,n)} \H_d(C^G_S(M))$ is finitely generated over
%$\Z[\sigma_1,\dots,\sigma_{k-2}]$. Condition (1) in 
%\cref{Prop:AlgStability} is satisfied by \cref{Thm:CGSHomStab}, and so
%$A_n = \bigoplus_{m\geq 0}A_{n,m}$ is finitely generated over
%$\Z[\sigma_1,\dots,\sigma_{k-1}]$.

%By the decomposition \eqref{eq:Hcal-decomp} combined with 
%\cref{Thm:CGSHomStab}, condition (1) in \cref{Prop:AlgStability} is
%satisfied with $A_n$ as above, $R = \Z[\sigma_1,\dots,\sigma_{k-1}]$, and
%$\sigma=\sigma_k$, and so $\coker(\sigma_k)$  is finitely generated as a
%$\Z[\sigma_1,\dots,\sigma_{k-1}]$-module. Here we use the fact that the underlying
%map of $\sigma_{G/H}$ induces the same map in homology as the iterated composite
%of the nonequivariant stabilization map (\cref{Lem:UnderlyingMap}). 

%%\EvaNote{Need to say that $\sigma_k$ after forgetting the action is the
%%classical stabilization map.}
%Lemma \ref{Lemma:Bredon-to-PhiG} (see also \cref{Rmk:ADWeak}) applied to $X = C(M)$ shows that
%$$ \coker(H^G_d(C(M);\underline{R}) \xrightarrow{\sigma_k}H_d^G(C(M);\underline{R})) $$
%is finitely generated over $\Z[\sigma_1,\dots,\sigma_{k-1}]$. Proposition
%\ref{Prop:AlgStability} applied to the sequence of groups $n\mapsto
%H^G_d(C_n(M))$ now shows that $\bigoplus_{n\geq 0} H^G_d(C_n(M)) = H^G_d(C(M))$
%is a finitely generated $\Z[\sigma_1,\dots,\sigma_k]$-module.
%\end{proof}

\subsection{Simplifying the hypotheses of Proposition \ref{Prop:BredonHomStab}}
We will show (Lemma \ref{lem:fg-homology}) that the finite generation hypothesis in Proposition
\ref{Prop:BredonHomStab} follows from the assumption that $M$
the interior of a compact $G$-manifold $W$. Such finite type hypotheses are
common in the literature on this subject; see for example \cite{RW13} for the
non-equivariant version of our result. We also translate the connectedness hypothesis of Proposition
\ref{Prop:BredonHomStab} into a condition
which involves $M^H$ instead of $M_{(H)}$ (see Lemma
\ref{lem:hypothesis-translation}). The results are combined in \cref{thm:B}.
\begin{lemma} \label{lem:hypothesis-translation}
Suppose $G$ is Dedekind or a $p$-group, and let $W = W_GH$ denote the Weyl
group. If $M^H/W$ is connected
%and $M^H$ is an open (i.e., noncompact) submanifold of $M$,
then $M_{(H)}/G$ (if nonempty) is
%open and
connected. Moreover,
$$ M_{(H)}/G \cong \Big(M^H \backslash \bigcup_{e < L \leq W} (M^H)^L\Big)/W. $$
\end{lemma}
\begin{proof}
Applying Lemma \ref{Lem:modW} to the $G$-manifold $M$ and the $W$-manifold
$M^H$, respectively, we have
\begin{align*}
M_{(H)}/G & \cong \Big(M^H \backslash \bigcup_{H<K\leq G} M^K\Big)\big/ W
\\ (M^H)_{(e)}/W & \cong \Big( M^H \backslash \bigcup_{e < L\leq W} (M^H)^L
\Big)\big/ W \cong \Big(M^H\backslash \bigcup_{H < K \leq N_G(H)} M^K\Big)\Big/W.
\end{align*}
Because of the group-theoretic conditions, $H < K \leq G$ implies $H < N_K(H)$, and
hence $M_{(H)}/G\cong (M^H)_{(e)}/W$.

If $M_{(H)}$ is nonempty, then $(e)$ is the principal orbit type of $M^H$ as a
$W$-manifold.
Then \cite[Theorem I.5.14]{tom-dieck-transformation-groups} applied to
$M^H$ says that $(M^H)_{(e)}/W \cong M_{(H)}/G$ is connected.
%and $(M^H)_{(e)}$ is dense in $M^H$. Since
%$M^H$ is noncompact, denseness implies that $(M^H)_{(e)}$ is noncompact, and
%hence $M_{(H)}/G = (M^H)_{(e)}/W$ is noncompact since the
%preimage of a compact set under a finite group quotient map is compact.
\end{proof}

The next lemma is used to prove Lemma \ref{lem:fg-homology}.

\begin{lemma}\label{lem:local-coeffs}
Suppose $M$ is $G$-manifold that is the interior of a compact $G$-manifold
with boundary. Let $H\leq G$. If $A$ is a $\pi_1(M^H)$-module that is finitely
generated as an abelian group, then the homology with local coefficients
$H_q(M^H;A)$ is a finitely generated abelian group for every $q$.
\end{lemma}
\begin{proof}
Let $Z$ be a compact $G$-manifold whose interior is $M$.
First, we observe that $M^H$ is the interior of $Z^H$;
this can be checked on Euclidean neighborhoods. Since $Z^H$ is a compact
manifold with boundary, it has a collar neighborhood, and this can be used to
produce a homotopy equivalence $M^H \hteq Z^H$. Since $\pi_1(M^H) = \pi_1(Z^H)$,
we may view $A$ as a sheaf of abelian
groups over $Z^H$, and consider an open cover $\mathscr{U}$ that trivializes
$A$. There is a refinement $\mathscr{U}'$ that is a good cover (see
\cite[Corollary 5.2]{bott-tu}) and since $Z^H$ is compact, we may take
$\mathscr{U}'$ to be finite. Since this is a good cover, homology with local
coefficients $H_q(Z^H;A)=H_q(M^H;A)$ agrees with \v{C}ech cohomology
$\check{H}_q(\mathscr{U}';A)$. The \v{C}ech complex
$$ \check{C}_q(\mathscr{U}';A) = \bigoplus_{(i_0,\dots,i_q)} A(U_{i_0}\ints
\dots \ints U_{i_q}) = \bigoplus_{(i_0,\dots,i_q)}A $$
is a finitely generated abelian group in each degree, and hence so is its homology.
\end{proof}

\begin{lemma} \label{lem:fg-homology}
Assume $G$ is abelian, and $M$ is the interior of a compact $G$-manifold.
Then $H_d(M_{(H)}/G; \Z)$ is finitely generated for all $d\in \Z,H\leq G$.
\end{lemma}
\begin{proof}
Let $W = W_GH$ be the Weyl group.
Lemma \ref{lem:hypothesis-translation} says that $M_{(H)}/G \cong X/W$ where
\begin{equation}\label{eq:bigcupX} X = M^H \backslash \bigcup_{e < L\leq W} (M^H)^L.
\end{equation}
By Lemma \ref{lem:homology-orbits-finite}, it suffices to show that
$H_*(X;\Z)$ is finitely generated in each
degree.

If $q:G\to W = G/H$ is the quotient, then $(M^H)^L = M^{q^{-1}(L)}$.
Abusing notation, we write $M^L = (M^H)^L$; also let
$U_L = M^H \backslash M^L$.
Since $M^L$ is a submanifold, there is a normal bundle $N(M^L)$ which is open,
and $N(M^L) \cup U_L \cong M^H$.
First we will show $H_d(U_L)$ is finite using a Mayer--Vietoris sequence.

The intersection term $N(M^L)\cap U_L$ is homotopy equivalent to the sphere bundle $S(N(M^L))$. There
is a Serre spectral sequence associated to the fibration $S^n\to S(N(M^L))\to
M^L$,
$$ E_2^{*,*} = H_*(M^L; H_*(S^n; \Z))\Rightarrow H_*(S(N(M^L));\Z) $$
where $\pi_1(M^L)$ may act nontrivially on $H_*(S^n;\Z)$. By Lemma
\ref{lem:local-coeffs}, the $E_2$ term is finitely generated in each degree, and
hence so is the target.

Now consider the Mayer--Vietoris sequence
$$ \dots\to H_{d+1}(M^H)\to H_d(S(N(M^L))) \to H_d(N(M^L)) \oplus H_d(U_L)\to
\dots. $$
The first term is finitely generated by Lemma \ref{lem:local-coeffs} and above we showed the second is
as well. Hence the third term is finitely generated; in particular so is $H_d(U_L)$.

Since $G$ is abelian, write $W = \prod_i \Z/p_i^{r_i}$.	In the union \eqref{eq:bigcupX}, it
suffices to restrict $L$ to products of elementary abelian subgroups
$\Z/p_i\cong \langle p_i^{r_i-1}\rangle\leq \Z/p_i^{r_i}$, since any nontrivial
subgroup $H\leq \Z/p_i^{r_i}$
contains this $\Z/p_i$. If $L\x L' \leq W$, we have $(M^H)^L \ints (M^H)^{L'} =
(M^H)^{L\x L'}$. More generally,
given $L = \prod_i L_i$ and $L' = \prod_i L'_i$, note
that $LL'\leq W$ is a product of all $L_i$ and $L'_i$ factors, with duplicates
removed. Thus we have
$$ U_L \cup U_{L'} = M^H \backslash ((M^H)^L \cap (M^H)^{L'}) = M^H \backslash
(M^H)^{LL'}. $$
We will show that 
$$X = \bigcap_\attop{L = L_1 \x \dots \x L_n\\L_i\text{ elem.ab.}} U_L $$
has finitely generated homology, by induction on the number of sets intersected.
The base case was shown above. Suppose any intersection of $n$ such sets has
finitely generated homology, and consider $U_{L_1}\ints \dots \ints U_{L_n}
\ints U_L$. There is a Mayer-Vietoris sequence
$$ \dots \to H_{d+1}((\cap_{i=1}^n U_{L_i})\cup U_L) \to H_d((\cap_{i=1}^n
U_{L_i})\cap U_L) \to H_d(\cap_{i=1}^n U_{L_i}) \oplus H_d(U_L) \to \dots. $$
The inductive hypothesis implies that the third term is finitely generated, so
in order to show the second term is finitely generated, it suffices to show the first
term is finitely generated. We have
$$ \Big(\bigcap_{i=1}^n U_{L_i}\Big) \cup U_L = \bigcap_{i=1}^n (U_{L_i} \cup
U_L) = \bigcap_{i=1}^n U_{L_i L} $$
and the inductive hypothesis shows this has finitely generated homology.
\end{proof}

\begin{theorem}\label{thm:B}
Let $G$ be an abelian group and let $M$ be a $G$-manifold that is the interior
of a compact $G$-manifold with boundary. Assume that for all $H\leq G$, $M$ is
$H$-stabilizable and $M^H$ is connected.
Then $H_d^H(C(M);\underline{\Z})$ is finitely generated over $P_G$ for all $H\leq G$.
\end{theorem}
\begin{proof}
Combine \cref{Prop:BredonHomStab}, \cref{lem:hypothesis-translation}, and
\cref{lem:fg-homology}.
\end{proof}

\begin{example} \label{ex:rho}
If $G$ is abelian and $V = n\rho_G$ is a sum of regular representations, the
hypotheses of Theorem \ref{thm:B} are satisfied. One could also check the
hypotheses of Proposition \ref{Prop:BredonHomStab} directly without using the
results of this subsection: by Lemma
\ref{Lem:modW} the space $V_{(H)}$ arises from the complement of hyperplanes in
$n\rho_G$. This complement is Alexander dual to a union of spheres, one for each
hyperplane.
\end{example}


\section{Generalizations and variants}\label{Sec:Gen}

Throughout this section, $G$ denotes a finite group in which every subgroup is normal. We have established stability for \emph{$\z$-graded} Bredon homology \emph{groups}. In this section, we discuss natural extensions of our results to
\begin{enumerate}
\item $RO(G)$-graded Bredon homology groups, and
\item $\z$-graded Bredon homology Mackey functors.
\end{enumerate}
Here, $RO(G)$ is the (group completion) of the ring of isomorphism classes of real orthogonal representations of $G$. 

\subsection{$RO(G)$-graded Bredon homological stability}

Let $V \in RO(G)$. We write $-V$ for the additive inverse of $V$ in $RO(G)$ and write $S^V$ for the representation sphere associated to $V$. 

\begin{definition}
Let $V \in RO(G)$. The $V$-th Bredon homology group of a $G$-space $X$ with coefficients in the Mackey functor $\uM$ is the group
$$\pi_0^G\left( S^{-V} \wedge H\uM \wedge X \right) \cong H_0^G( S^{-V} \wedge X; \uM).$$
\end{definition}

The following is well-known:

\begin{lemma}\label{Lem:RepSpheres}
If $V \in RO(G)$, then we have 
$$(S^V)^e = \Phi^e(S^V) = S^{|V|}, \quad  \Phi^G(S^V) = S^{|V^G|}.$$
\end{lemma}

Recall that $P_G = \Z[(\sigma_{G/H})_* : (H)\leq G]$.
\begin{theorem}\label{Thm:ROG}
Let $V \in RO(G)$ and let $M$ be a $G$-manifold satisfying the hypotheses of
\cref{thm:B}. Then the $P_G$-module $H_V^G(C(M);\mfz)$ is finitely generated. 
\end{theorem}

\begin{proof}
Applying \cref{Lemma:Bredon-to-PhiG} with $X = S^{-V} \wedge C(M)$ and $\cc$ the Serre class of finitely generated $P_G$-modules, it suffices to show that $\pi_q \Phi^K(H\mfz \wedge S^{-V} \wedge C(M))$ is in $\cc$ for all $K \leq G$ and all $q \in \z$. \cref{Lem:RepSpheres} implies that
$$\pi_q \Phi^K \left(H\mfz \wedge S^{-V} \wedge C(M) \right) \cong \pi_{q+|V^K|} \Phi^K \left(H\mfz \wedge C(M) \right),$$
and this is finitely generated by combining the same ingredients as the proof of
Proposition \ref{Prop:BredonHomStab} and Theorem \ref{thm:B}.
\end{proof}


\subsection{Mackey functor stability}

If $L \leq G$ is a subgroup, we may view any $G$-manifold $M$ as an $L$-manifold by restriction. If the $G$-manifold $M$ satisfies the hypotheses of \cref{MT:Stability}, then its underlying $L$-manifold clearly also satisfies the conditions of \cref{MT:Stability} (with $L$ in place of $G$) and we obtain the following:

\begin{corollary}
Suppose $M$ satisfies the hypotheses of \cref{thm:B}.
Then $H_d^L(C(M);\mfz)$ is finitely generated over $P_L$ for all integers $d \geq 0$.  
\end{corollary}

Recall that Bredon homology is actually a Mackey functor, i.e., an additive functor from the Burnside category of $G$ to abelian groups. The previous corollary implies that the value of the Mackey functor
$$H_d^{(-)}(C(M); \mfz)$$
stabilizes when evaluated on any object in the Burnside category (i.e., any finite $G$-set). We can encode this levelwise stability, and more, with the following notion:

\begin{definition}
A \emph{module} over a Green functor $\uR$ is a Mackey functor $\uM$ equipped with an action map $\uR \boxtimes \uM \to \uM$ satisfying the usual associativity and unitality conditions. 

An $\uR$-module $\uM$ is \emph{finitely generated} if there is a surjective $\uR$-linear map $\uR\{x_T\} \to \uM$, where $\uR\{x_T\}$ is the free $\uR$-module on a generator $x_T$ with $T$ a finite $G$-set. 
\end{definition}

According to \cite[Lem. 2.17]{Shu10}, a map $\uR \boxtimes \uM \to \uM$ determines, and is determined by, its  \emph{Dress pairing}: a collection of maps $\uR(S) \otimes \uM(S) \to \uM(S)$ for $S$ in the Burnside category of $G$ satisfying certain compatibility axioms.  Using Dress pairings, we can relate levelwise finite generation to Mackey finite generation:

\begin{lemma}\label{Lem:FGLevelwise}
An $\uR$-module $\uM$ is finitely generated if and only if $\uM(S)$ is finitely generated as a $\uR(S)$-module for each finite $G$-set $S$.
\end{lemma}

\begin{proof}
Surjectivity for Mackey functors is checked levelwise. 
\end{proof}

To express stability in terms of Mackey functors, we need a Green functor to encode the action of the equivariant stabilization maps for each subgroup of $G$. This is accomplished through the following:

\begin{definition}
For each finite group $G$, let $\uP_G$ be the constant Green functor on $P_G=\z[(\sigma_{G/H})_*: (H) \leq G]$. (\emph{A priori}, $\uP_G$ is only a Mackey functor. It is a Green functor since it is the fixed point functor of a commutative ring with trivial action.) We define an action of $\uP_G$ on $H_*^{(-)}(C(M); \mfz)$ via Dress pairing: the map
$$\uP_G(G/H) \otimes H_*^H(C(M);\mfz) \cong P_G\otimes H^H_*(C(M);\mfz) \to H_*^H(C(M);\mfz)$$
comes from the restriction map $\res^G_H:P_G\to P_H$ as defined in Section \ref{SS:BredonStab}.
\end{definition}

In particular, if $H_*^H(C(M);\mfz)$ is finitely generated over $\uP_H(H/H)$ for every subgroup $H \leq G$, and if $\uP_H(H/H)$ is finitely generated as a $\uP_G(G/H)$-module via the restriction map (e.g., if every subgroup of $G$ is normal), then $H_*^H(C(M);\mfz)$ is finitely generated as a $\uP_G(G/H)$-module. Combined with \cref{Lem:FGLevelwise}, this proves:

\begin{theorem}\label{Thm:MF}
Let $M$ be as in \cref{thm:B}. Then $H_d^{(-)}(C(M); \mfz)$ is finitely generated over $\uP_G$ for all integers $d \geq 0$.  
\end{theorem}


\bibliographystyle{alpha}
\bibliography{master}

\end{document}
